% TeX root=../main.tex
\setlength{\textfloatsep}{5pt}

\section{Substantivos}

\subsection{Categorias}

As categorias nominai são gênero (masculino, feminino, neutro), número
(singular, dual e plural) e caso (nominativo, acusativo, genitivo, dativo,
instrumental, locativo, vocativo).

\emph{Caso}: o eslavo antigo preserva a situação herdada do indo-europeu, a
única modificação estando limitada a um caso de sincretismo no sistema de casos,
entre o genitivo e o ablativo (os quais já haviam se confundido parcialmente em
indo-europeu): nos temas em \emph{-o-}, é a forma antiga do ablativo que, na
sincronia, é a forma de genitivo, o que é uma inovação de data balto-eslávica;
nas demais declinações, o genitivo continua a forma etimológica do genitivo.
Os valores principais dos casos são os seguintes:

% TeX root=../main.tex

\begin{table}[!htb]
	\caption{Valores principais dos casos}\label{tab:valor_dos_casos}
	\begin{center}
		\begin{tabular}[c]{lll}
			\toprule
			                    & Gramatical                                         & Concreto               \\
			\midrule
			\emph{Nominativo}   & sujeito                                            &                        \\
			\emph{Acusativo}    & objeto                                             & direção                \\
			\emph{Genitivo}     & complemento nominal                                & origem                 \\
			\emph{Dativo}       & objeto indireto                                    & atribuição, destinação \\
			\emph{Instrumental} & agente da passiva                                  & meio, maneira          \\
			\emph{Locativo}     &                                                    & lugar                  \\
			\midrule
			\emph{Vocativo}     & \multicolumn{2}{l}{interlocutor (fora da sintaxe)}                          \\
			\bottomrule
		\end{tabular}
	\end{center}
\end{table}


\emph{Número}: certos substantivos são \emph{pluralia tantum} (\gloss{neut.pl.}
\cirilocsex{vrata} `portas', \gloss{masc.pl.} \cirilocsex{l'judije,l'udije}
`gentes'), enquanto outros são \emph{singularia tantum} (os nomes próprios, os
coletivos como os \gloss{fem.sg.} \cirilocsex{CEdI} `gente',
\cirilocsex{bratija} `conjunto de irmãos').

\emph{Gênero}: o neutro é caracterizado pela ausência de distinção entre
nominativo e acusativo, mas essa regra conhece uma exceção no caso dos
particípios; sua flexão não se distingue daquela do masculino se não no
nominativo, acusativo e vocativo, como é a regra nas línguas indo-europeias.
A atribuição deste ou daquele gênero a um substantivo não obedece nenhuma
consideração semântica, mas a critérios morfológicos (gênero gramatical): neste
sentido, gênero é arbitrário.

O \glsxtrlong{eslavocomum}, no entanto, inovou ao criar uma distinção suplementar, que
obedece a um critério semântico, no interior do gênero masculino, entre os
animados e inanimados: o inanimado é caracterizado por uma forma de acusativo
idêntica à do nominativo, o animado por uma forma de acusativo idêntica à do
genitivo, apenas no singular.
Nos demais casos em que não o acusativo, animados e inanimados possuem as mesmas
formas.
Há assim para os animados uma neutralização da oposição entre caso acusativo e
genitivo e, para os inanimados, da oposição entre o caso nominativo e acusativo.
A concordância produz uma forma acusativa animada (idêntica aos genitivos) para
os determinantes de um animado.
Nos substantivos, o subgênero animado não existe em eslavo antigo salvo nos
masculinos que designam pessoas (os nomes de animais são tratados como
inanimados) e somente no singular.
Suas extensão é maior nos pronomes, onde existe também para o plural.
Nas línguas eslavas modernas, o subgênero animado sofreu uma certa extensão e se
encontra também no plural, e mesmo ao feminino em eslavo oriental.
É uma forma de marcação diferenciada de objeto (o objeto animado tem uma marca
específica que o inanimado não tem).

Historicamente, esse fato se explica pela necessidade de distinguir o nominativo
e o acusativo dos temas em -\ocsex{o}- (antiga classe temática indo-europeia),
a mudança fonológica tornara homófonos os dois casos no singulas
(\glsxtrshort{eslavoantigo}~\gloss{nom.sg.}~\ocsex{-U},
\mbox{\gloss{acu.sg.}~\ocsex{-U}}): a necessidade de distinguir agente e paciente nos
animados de uma língua na qual a ordem de palavras é livre conduziu a língua a
inovar, passando a empregar na função do acusativo a forma genitiva.
Este emprego secundário foi possibilitado por uma característica sintática do
\glsxtrlong{eslavocomum}: o complemento de objeto de um verbo negativo é no genitivo e não
no acusativo (genitivo de negação, \emph{cf.}~\ref{sec:negacao}); em uma oração
negativa, a oposição entre acusativo e genitivo era assim neutralizada, e é uma
forma do genitivo que compre a função do acusativo.
O uso da forma genitiva em função de acusativo para os animados é uma extensão
dessa neutralização de oposição entre os dois casos: há simplesmente uma troca
de critério, o critério determinate na neutralização não sendo mais sintático
(negação), mas semântico (traço animado).

Em \glsxtrlong{eslavoantigo}, todos os substantivos animados de tema em
-\ocsex{o}- não possuem obrigatoriamente a forma de \gloss{gen.-acu.}. De fato,
mais critérios determina a aparição da forma de \gloss{gen.-acu.}, além do
critério fundamental (animacidade).
Eles são: a posição do substantivo na hierarquia dos animados (um substantivo próprio é
hierarquicamente superior a um substantivo comum que designa uma pessoa e o
\gloss{gen.-acu.} será mais sistemático entre substantivos próprios), a
determinação (a forma de animado é empregada para um substantivo determinado,
\ocsex{viZdõ ClovEka} `eu vejo o homem', mas não para um substantivo
indeterminado, \ocsex{viZdõ ClovEkU} `eu vejo um homem'), porque a determinação
implica numa especificação superior, tal como o substantivo próprio é mais
especificado que o comum), a sintaxe (o \gloss{gen.-acu.} é menos frequente nas
construções preposicionais que o acusativo do que no complemento de objeto:
\ocsex{inogo} \textsubscript{\gloss{gen.-acu.}}\ocsex{posUla} `ele envia um
outro' \passageme{M125}, mas \ocsex{posagnetU za
	inU}\textsubscript{\gloss{acu.}} `(se) ela esposa um outro' \passageme{M1012}).

Em \glsxtrlong{eslavoantigo}, observa-se a tendência à extensão do subgênero
animado para além do acusativo: o \gloss{dat.sg.} dos mesmos temas em
-\ocsex{o}-, cuja desiência é -\ocsex{u}, conhece um alomorfo -\ocsex{ovi},
emprestado dos temas em -\ocsex{u}-, e que não se encontra se não em
substantivos animados, a língua introduz assim uma diferença entre o dativo
animado (-\ocsex{ovi}) e o dativo inanimado (-\ocsex{u}).
Essa distribuição está longe de ser sistematizada em \glsxtrlong{eslavoantigo},
onde ela constitui apenas uma tendência.
Observa-se da mesma sorte uma extensão da forma de \gloss{gen.-acu.} aos
substantivos animado de outros tipos flexionais, como \cirilocsex{synU} `filho'
(tema em -\ocsex{u}-) ou \cirilocsex{gospodI} `senhor' (tema em -\ocsex{i}-),
bem como a femininos (\cirilocsex{mati} `mãe').


\subsection{Paradigmas}

Os paradigmas flexionais são cinco em número:
temas em -\ocsex{o}- (com um subtipo em -\ocsex{jo}-), herdados da flexão
temática do indo-europeu;
temas em -\ocsex{a}- (com um subtipo em -\ocsex{ja}-), provenientes dos temas em
*-\pietrans{eh2},
temas em -\ocsex{i}-, temas em -\ocsex{u}- e temas em consoante.
Os três primeiros continuam vivos nas línguas eslavas modernas.
Em \glsxtrlong{eslavoantigo}, os três primeiros ainda são produtivos e os dois
últimos são recessivos: eles são substituídos pelos tipos produtivos e possuem
por isso uma dupla flexão; eles não acolhem empréstimos salvo exceções.
Há uma correlação entre o paradigma flexional e o gênero da palavra: os temas em
-\ocsex{o}- são masculinos ou neutros, os temas em -\ocsex{a}- são femininos,
com alguns exemplos raros masculinos, os temas em -\ocsex{u}- são masculinos.
Isso conserva a situação indo-europeia.
Para os temas em -\ocsex{i}-, eles são sobretudo femininos, mas contém também
alguns masculinos que tem a tendência a passar para temas em -\ocsex{jo}-, sendo
a tendência de eliminação dos masculinos em -\ocsex{i}- já do
\glsxtrlong{eslavocomum}.
Os temas consonantais, por fim, possuem substantivos dos três gêneros e se
dividem em outros subtipos.

\subsubsection{Temas vocálicos}\label{sec:temasvocalicos}

Os temas em -\emph{o}- (\autoref{tab:temasemo}) e em -\emph{a-}
(\autoref{tab:temasema}) se dividem
em dois subtipos, chamados ``tipo duro'' (temas em -\emph{o}- e
-\emph{a}-) e ``tipo mole'' (temas em -\emph{jo}- e -\emph{ja}-). Essa
cisão é uma consequência da harmonia intrassilábica do eslavo comum (\emph{cf.}~\ref{sec:estruturasilabica}): à medida em que uma vogal posterior tem seu ponto de
articulação avançado até o seu correspondente palatal após /j/ (e, por
conseguinte, em sincronia, igualmente após uma consoante palatalizada
resultante de uma sequência C + /j/), a vogal temática *\emph{o}
torna-se *\emph{e} nos temas em -\emph{jo}- (tipo mole), enquanto ela
continua \emph{o} nos temas em \emph{-o-} (tipo duro). A evolução
fonética subsequente acentuou a divergência entre as variantes: assim,
no \gloss{loc.masc.sg.} *-\emph{oi̯} > -\emph{ě} (tipo duro), mas
*\emph{-i̯oi̯} > *\emph{i̯ei̯} > -\emph{{(j)}i}
(tipo mole), no \gloss{acu.masc.pl}., \emph{ons} > *-\emph{uns}
> *-\emph{ūs} > \emph{-y} (tipo duro), mas
*-\emph{{(i̯)}ons} > *-\emph{{(i̯)}ens} >
\emph{{(j)}ę} (tipo mole). O mesmo processo vale para os temas em
-\emph{a}- e os temas em -\emph{ja}-, largamente paralelos aos temas em
-\emph{o}- e -\emph{jo}-. Encontram-se, portanto, dois subtipos e duas
séries de desinências: embora a relação entre os dois seja transparente
para algumas formas casuais, isso não ocorre com todas (assim, no
\gloss{acu.pl.masc.} e \gloss{fem.} em que a desinência de tipo duro -\emph{y} foi
desnasalizada enquanto a de tipo mole -\emph{ę} continua uma vogal
nasal).

\subparagraph{Alternâncias:} no interior do tipo duro (temas em -\emph{o}- e em
-\emph{a-}), os temas em velar são caracterizados por uma alternância
morfofonológica da consoante final \mbox{\emph{k} $\sim$ \emph{c}
$\sim$ \emph{č}}, \mbox{\emph{g} $\sim$ \emph{dz} $\sim$ \emph{ž}},
\mbox{\emph{x} $\sim$ \emph{s} $\sim$ \emph{š}}, consequência de
palatalizações sucessivas (primeira palatalização para as formas em
\emph{č}, \emph{ž}, \emph{š} que aparecem no \gloss{voc.masc.sg.}; segunda
palatalização para as formas em \emph{c}, \emph{dz}, (>
\emph{z}), \emph{s} que aparecem no masculino no \gloss{loc.sg.} e \gloss{pl.} e
\gloss{nom.pl.},
no neutro no \gloss{loc.sg.} e \gloss{pl.} e \gloss{nom.acu.dual.}, e no
feminino no \gloss{dat.loc.sg.}
e \gloss{nom.acu.dual}). Nos temas em -\emph{sk}-, a forma palatalizada
-\emph{scě}, \emph{-sci} é simplificada em \emph{-stě}, -\emph{sti} em
certos manuscritos (\glsxtrshort{Supr}). O substantivo \cirilocsex{vlUxvU}
``mago'' apresenta a alternância regular nos temas em velar mesmo que o
[x] não esteja em posição final do tema: \gloss{nom.pl.} \cirilocsex{vlUsvi}
(\passageme{Mt21} \glsxtrshort{Assemanius}), o que resulta do fato de que os
grupos \emph{*kv}, *\emph{gv}, *\emph{xv} são afetadas pela segunda
palatalização e se tornam \emph{cv}, \emph{(d)zv}, \emph{sv}.
No tipo
mole, os substantivos em \cirilocsex{-IsI} (temas em -\emph{o}-)
apresentam uma alternância \emph{c $\sim$ č} no \gloss{voc.sg.}:
\gloss{nom.sg.} \cirilocsex{otIcI} ``pai'', \gloss{gen.sg.} \cirilocsex{otIca},
\gloss{voc.sg.} \cirilocsex{otICe}.

\subparagraph{Sincretismo e homofonias:} no tipo feminino em -\emph{a}-, o
sincretismo do \gloss{dat.loc.sg.} e do \gloss{nom.acu.pl.} é sistemático: o primeiro é
herdado do indo-europeu (*\emph{-āi} < Dat. *\emph{-āei} e
loc. \emph{*-āi}), o segundo é de época do eslavo comum. O tipo em
-\emph{o}- masc. apresenta ademais uma homofonia entre \gloss{nom.acu.sg.} e
\gloss{gen.pl.} (\emph{-ъ} no tipo duro, -\emph{ь} no tipo mole), entre \gloss{gen.sg.} e
\gloss{nom.acu.du.}, e, unicamente no tipo duro, entre \gloss{acu.pl.masc.} e
\gloss{inst.pl.}
(-\emph{y}). No tipo em -\emph{o}-, a oposição entre tipo duro e tipo
mole é neutralizada no \gloss{nom.pl.masc.} (-\emph{i-}). A oposição de gênero
entre os temas em -\emph{o}- e em -\emph{a}- é neutralizada nos
seguintes casos: \gloss{loc.sg.} (três gêneros), \gloss{gen.loc.dual} (três gêneros),
\gloss{gen.pl.} (três gêneros), \gloss{acu.pl.} (somente \gloss{masc.} e
\gloss{fem.}, os neutros têm uma desinência específica), \gloss{nom.acu.dual}
(\gloss{fem.} e \gloss{nt.} somente). A adoção pelo tipo em -\emph{o}- de certas
desinências do tipo em -\emph{u}- consiste em parte em uma resposta a essa
subespecificação (adiante).
\clearpage

% TeX root=../main.tex
\begin{table}[!htb]
	\begin{adjustwidth}{0cm}{-2cm}
		\caption[Temas em -\ocsex{o}-]{Temas em -\ocsex{o}- (Paradigmas: \gloss{masc.}~\ocsex{gradU} `cidade', \ocsex{ClovEkU} `homem', \ocsex{kon'I} `cavalo';
			\gloss{neut.}~\ocsex{mEsto} `lugar', \ocsex{mor'e} `mar')}\label{tab:temasemo}
		\begin{center}
			\begin{tabular}[c]{llllllll}
				\toprule
				                   &                   & \multicolumn{2}{c}{Duro inanimado} & \multicolumn{2}{c}{Duro animado} &
				\multicolumn{2}{c}{Mole inanimado}                                                                                                                                                               \\
				\midrule
				\emph{sg.}         & \emph{nom.}       & \cirilex{gradU}                    & \ocsex{grad-U}                   & \cirilex{ClovEkU}   & \ocsex{ClovEk-U}   & \cirilex{kon'I}   & \ocsex{kon'-I}   \\
				\rowcolor{gray!10} & \emph{acu.}       & \cirilex{gradU}                    & \ocsex{grad-U}                   & \cirilex{ClovEka}   & \ocsex{ClovEk-a}   & \cirilex{kon'I}   & \ocsex{kon'-I}   \\
				                   & \emph{n.-a.neut.} & \cirilex{mEsto}                    & \ocsex{mEsto}                    &                     &                    & \cirilex{mor'e}   & \ocsex{mor'-e}   \\
				\rowcolor{gray!10} & \emph{gen.}       & \cirilex{grada}                    & \ocsex{grad-a}                   & \cirilex{ClovEka}   & \ocsex{ClovEk-a}   & \cirilex{kon'a}   & \ocsex{kon'-a}   \\
				                   & \emph{dat.}       & \cirilex{gradu}                    & \ocsex{grad-u}                   & \cirilex{ClovEku}   & \ocsex{ClovEk-u}   & \cirilex{kon'u}   & \ocsex{kon'-u}   \\
				\rowcolor{gray!10} & \emph{inst.}      & \cirilex{gradomI}                  & \ocsex{grad-omI}                 & \cirilex{ClovEkomI} & \ocsex{ClovEk-omI} & \cirilex{kon'emI} & \ocsex{kon'-emI} \\
				                   & \emph{loc.}       & \cirilex{gradE}                    & \ocsex{grad-E}                   & \cirilex{ClovEkE}   & \ocsex{ClovEk-E}   & \cirilex{kon'i}   & \ocsex{kon'-i}   \\
				\rowcolor{gray!10} & \emph{voc.}       & \cirilex{grade}                    & \ocsex{grad-e}                   & \cirilex{ClovEke}   & \ocsex{ClovEk-e}   & \cirilex{kon'u}   & \ocsex{kon'-u}   \\
				\midrule
				\emph{du.}         & \emph{nom.acu}    & \cirilex{grada}                    & \ocsex{grad-a}                   & \cirilex{ClovEka}   & \ocsex{ClovEk-a}   & \cirilex{kon'a}   & \ocsex{kon'-a}   \\
				\rowcolor{gray!10} & \emph{n.-a.neut.} & \cirilex{mEstE}                    & \ocsex{mEst-E}                   &                     &                    & \cirilex{mor'i}   & \ocsex{mor'-i}   \\
				                   & \emph{gen.loc.}   & \cirilex{gradu}                    & \ocsex{grad-u}                   & \cirilex{ClovEku}   & \ocsex{ClovEk-u}   & \cirilex{kon'u}   & \ocsex{kon'-u}   \\
				\rowcolor{gray!10} & \emph{dat.inst.}  & \cirilex{gradoma}                  & \ocsex{grad-oma}                 & \cirilex{ClovEkoma} & \ocsex{ClovEk-oma} & \cirilex{kon'ema} & \ocsex{kon'-ema} \\
				\midrule
				\emph{pl.}         & \emph{nom.}       & \cirilex{gradi}                    & \ocsex{grad-i}                   & \cirilex{ClovEki}   & \ocsex{ClovEk-i}   & \cirilex{kon'i}   & \ocsex{kon'-i}   \\
				\rowcolor{gray!10} & \emph{acu.}       & \cirilex{grady}                    & \ocsex{grad-y}                   & \cirilex{ClovEky}   & \ocsex{ClovEk-y}   & \cirilex{kon'ẽ}   & \ocsex{kon'-ẽ}   \\
				                   & \emph{n.-a.neut.} & \cirilex{mEsta}                    & \ocsex{mEst-a}                   &                     &                    & \cirilex{mor'ja}  & \ocsex{mor'-a}   \\
				\rowcolor{gray!10} & \emph{gen.}       & \cirilex{gradU}                    & \ocsex{grad-U}                   & \cirilex{ClovEkU}   & \ocsex{ClovEk-U}   & \cirilex{kon'I}   & \ocsex{kon'-I}   \\
				                   & \emph{dat.}       & \cirilex{gradomU}                  & \ocsex{grad-omU}                 & \cirilex{ClovEkomU} & \ocsex{ClovEk-omU} & \cirilex{kon'emU} & \ocsex{kon'-emU} \\
				\rowcolor{gray!10} & \emph{inst.}      & \cirilex{grady}                    & \ocsex{grad-y}                   & \cirilex{ClovEky}   & \ocsex{ClovEk-y}   & \cirilex{kon'i}   & \ocsex{kon'-i}   \\
				                   & \emph{loc.}       & \cirilex{gradExU}                  & \ocsex{grad-ExU}                 & \cirilex{ClovEkExU} & \ocsex{ClovEk-ExU} & \cirilex{kon'ixU} & \ocsex{kon'-ixU} \\
				\bottomrule
			\end{tabular}
		\end{center}
	\end{adjustwidth}
\end{table}


Os temas em -\emph{o}- masculinos apresentam também uma série de
desinências diferentes, emprestadas aos temas em -\emph{u}- (\autoref{tab:temasemu});
essas desinências concorrentes são reservadas aos masculinos e não são
encontradas nos neutros.
\clearpage


% TeX root=../main.tex

\begin{table}[!htb]
	\begin{adjustwidth}{0cm}{0cm}
		\caption[Temas em -\ocsex{a}-]{Temas em -\ocsex{a}-\\ (Paradigmas:
			\ocsex{Zena} `mulher', \ocsex{rEka} `rio', \ocsex{duSa} `alma')}\label{tab:temasema}
		\begin{center}
			\begin{tabular}[c]{llllllll}
				\toprule
				                   &                  & \multicolumn{2}{c}{Duro} & \multicolumn{2}{c}{Duro} &
				\multicolumn{2}{c}{Mole}                                                                                                                                              \\
				\midrule
				\emph{sg.}         & \emph{nom.}      & \cirilex{Zena}           & \ocsex{Zen-a}            & \cirilex{rEka}   & \ocsex{rEk-a}   & \cirilex{duSa}   & \ocsex{duS-a}   \\
				\rowcolor{gray!10} & \emph{acu.}      & \cirilex{Zenõ}           & \ocsex{Zen-õ}            & \cirilex{rEkõ}   & \ocsex{rEk-õ}   & \cirilex{duSõ}   & \ocsex{duS-õ}   \\
				                   & \emph{gen.}      & \cirilex{Zeny}           & \ocsex{Zen-y}            & \cirilex{rEky}   & \ocsex{rEk-y}   & \cirilex{duSẽ}   & \ocsex{duS-ẽ}   \\
				\rowcolor{gray!10} & \emph{dat.}      & \cirilex{ZenE}           & \ocsex{Zen-E}            & \cirilex{rEcE}   & \ocsex{rEc-E}   & \cirilex{duSi}   & \ocsex{duS-i}   \\
				                   & \emph{inst.}     & \cirilex{Zenojõ}         & \ocsex{Zen-ojõ}          & \cirilex{rEkojõ} & \ocsex{rEk-ojõ} & \cirilex{duSejõ} & \ocsex{duS-ejõ} \\
				\rowcolor{gray!10} & \emph{loc.}      & \cirilex{ZenE}           & \ocsex{Zen-E}            & \cirilex{rEcE}   & \ocsex{rEc-E}   & \cirilex{duSi}   & \ocsex{duS-i}   \\
				                   & \emph{voc.}      & \cirilex{Zeno}           & \ocsex{Zen-o}            & \cirilex{rEko}   & \ocsex{rEk-o}   & \cirilex{duSe}   & \ocsex{duS-e}   \\
				\midrule
				\emph{du.}         & \emph{nom.acu}   & \cirilex{ZenE}           & \ocsex{Zen-E}            & \cirilex{rEcE}   & \ocsex{rEc-E}   & \cirilex{duSi}   & \ocsex{duS-i}   \\
				\rowcolor{gray!10} & \emph{gen.loc.}  & \cirilex{Zenu}           & \ocsex{Zen-u}            & \cirilex{rEku}   & \ocsex{rEk-u}   & \cirilex{duSu}   & \ocsex{duS-u}   \\
				                   & \emph{dat.inst.} & \cirilex{Zenama}         & \ocsex{Zen-ama}          & \cirilex{rEkama} & \ocsex{rEk-ama} & \cirilex{duSama} & \ocsex{duS-ama} \\
				\midrule
				\emph{pl.}         & \emph{nom.}      & \cirilex{Zeny}           & \ocsex{Zen-y}            & \cirilex{rEky}   & \ocsex{rEk-y}   & \cirilex{duSẽ}   & \ocsex{duS-ẽ}   \\
				                   & \emph{acu.}      & \cirilex{Zeny}           & \ocsex{Zen-y}            & \cirilex{rEky}   & \ocsex{rEk-y}   & \cirilex{duSẽ}   & \ocsex{duS-ẽ}   \\
				\rowcolor{gray!10} & \emph{gen.}      & \cirilex{ZenU}           & \ocsex{Zen-U}            & \cirilex{rEkU}   & \ocsex{rEk-U}   & \cirilex{duSI}   & \ocsex{duS-I}   \\
				                   & \emph{dat.}      & \cirilex{ZenamU}         & \ocsex{Zen-amU}          & \cirilex{rEkamU} & \ocsex{rEk-amU} & \cirilex{duSamU} & \ocsex{duS-amU} \\
				\rowcolor{gray!10} & \emph{inst.}     & \cirilex{Zenami}         & \ocsex{Zen-ami}          & \cirilex{rEky}   & \ocsex{rEk-ami} & \cirilex{duSami} & \ocsex{duS-ami} \\
				                   & \emph{loc.}      & \cirilex{ZenaxU}         & \ocsex{Zen-axU}          & \cirilex{rEkaxU} & \ocsex{rEk-axU} & \cirilex{duSixU} & \ocsex{duS-axU} \\
				\bottomrule
			\end{tabular}
		\end{center}
	\end{adjustwidth}
\end{table}


O tipo em -\emph{a}- mole contém um subtipo no \gloss{nom.sg.} em
\cirilocsex{-i}, que ocorre em poucos substantivos, mais em femininos de
particípios e do comparativo (\emph{cf.}~\ref{sec:comparativo} e
\ref{sec:participios}).
Para esse subtipo, a única forma divergente é a do \gloss{nom.sg.} em -\emph{i}
(\gloss{fem.} \cirilocsex{bogyn'i} ``deusa'', \gloss{masc.} \cirilocsex{sõdiji}
``juiz'', \gloss{acu.sg.} \cirilocsex{bogyn'jõ,bogyn'õ}, \cirilocsex{sõdijõ},
\gloss{gen.sg.} \cirilocsex{bogyn'ẽ}, \cirilocsex{sõdijẽ} etc.), mas ela concorre com uma
desinência nova \cirilocsex{ja}, analógica do tipo dominante.
Notar-se-á
que entre os substantivos em \cirilocsex{iji} aparecem alguns
masculinos, raros de modo geral no tipo em -\emph{a}- (exceto em
compostos, \emph{cf.}~\ref{sec:composicao}).

O paradigma dos temas em -\emph{u}- (\autoref{tab:temasemu}) caracterizada pela
homofonia entre gen.sg. e loc.sg., não é mais produtivo em AEE e está em
vias de ser absorvido pelo tipo em -\emph{o}- (os dois tipos têm a mesma
desinência -\emph{ъ} no nom.acu.sg., o que facilita a passagem de um ao
outro, e o tipo em -\emph{u}- contém apenas masculinos).
\clearpage


% TeX root=../main.tex

\begin{table}
  \caption[Temas em \ocsex{-u-}]{Temas em \ocsex{-u-} (Paradigma: \ocsex{synU} `filho')}\label{tab:temasemu}
	\begin{center}
		\begin{tabular}[c]{lllllll}
			\toprule
			                                & \multicolumn{2}{c}{Singular} & \multicolumn{2}{c}{Dual} & \multicolumn{2}{c}{Plural}                                                        \\
			\midrule
			\rowcolor{gray!10} \emph{nom.}  & \cirilex{synU}               & \ocsex{syn-U}            & \cirilex{syny}             & \ocsex{syn-y}   & \cirilex{synove} & \ocsex{syn-ove} \\
			\emph{acu.}                     & \cirilex{synU}               & \ocsex{syn-U}            & \cirilex{syny}             & \ocsex{syn-y}   & \cirilex{syny}   & \ocsex{syn-y}   \\
			\rowcolor{gray!10} \emph{gen.}  & \cirilex{synu}               & \ocsex{syn-u}            & \cirilex{synovu}           & \ocsex{syn-ovu} & \cirilex{synovU} & \ocsex{syn-ovU} \\
			\emph{dat.}                     & \cirilex{synovi}             & \ocsex{syn-ovi}          & \cirilex{synUma}           & \ocsex{syn-Uma} & \cirilex{synomU} & \ocsex{syn-omU} \\
			\rowcolor{gray!10} \emph{inst.} & \cirilex{synUmI}             & \ocsex{syn-UmI}          & \cirilex{synUma}           & \ocsex{syn-Uma} & \cirilex{synUma} & \ocsex{syn-Umi} \\
			\emph{lac.}                     & \cirilex{synu}               & \ocsex{syn-u}            & \cirilex{synovu}           & \ocsex{syn-ovu} & \cirilex{synoxU} & \ocsex{syn-oxU} \\
			\rowcolor{gray!10} \emph{voc.}  & \cirilex{synu}               & \ocsex{syn-u}            &                            &                 &                  &                 \\
			\bottomrule
		\end{tabular}
	\end{center}
\end{table}


Consequentemente, encontram-se desinências do tipo em -\emph{o}- nos temas em
-\emph{u}- etimológicos (\gloss{gen.sg.} \cirilocsex{syna} em vez de
\cirilocsex{synu}, \gloss{gen.-acu.} animado \cirilocsex{syna} em vez de
\cirilocsex{synU}, \gloss{loc.sg.} \cirilocsex{synE} em vez de
\cirilocsex{synu}, \gloss{inst.sg.} \cirilocsex{synomI} em vez de
\cirilocsex{synUmI}), e assim, inversamente, as desinências do tipo em
-\emph{u}- nos temas em -\emph{o}-.
O tipo em -\emph{o}- notadamente emprestou o \gloss{gen.pl.} em
\emph{-ovъ} dos temas em -\emph{u}- (\cirilocsex{grExovU} ``dos
pecados'' em lugar de \cirilocsex{grExU}), o que permite, por um lado,
instaurar uma oposição de gênero entre masc. (-\emph{ovъ}) e \emph{fem.nt.}
(-\emph{ъ}) e, por outro, restaurar uma oposição entre o
\gloss{nom.acu.sg.masc.} e \gloss{gen.pl.masc.}, homófonos no tipo em -\emph{o}-
(-\emph{ъ}); essa tendência, que começa a se manifestar em AEE, terá um
grande desenvolvimento nas línguas eslavas modernas. O tipo em
-\emph{o}- também emprestou o \gloss{dat.sg.} em -\emph{ovi}, que se estende às
custas de -\emph{u}, desinência etimológica do tipo em \emph{-o},
seguindo uma repartição semântica: o \gloss{dat.} em -\emph{ovi} é reservado aos
substantivos animados (\cirilocsex{bogovi} ``a Deus'', e para o tipo
mole  \cirilocsex{mõZevi} ``ao homem''), os inanimados conservam
-\emph{u}; o subgênero animado tende, assim, a se estender para além de
seu contexto de aparição que é o \gloss{acu.sg.} e ocorre também em uma outra
forma casual. As desinências de gen.sg. -\emph{u}, loc.sg. -\emph{u},
\gloss{nom.pl.} \emph{-ove} dos temas em -\emph{u}- são igualmente atestados
ocasionalmente no tipo em -\emph{o}- (\gloss{nom.pl.} \cirilocsex{darove},
\gloss{loc.sg.} \cirilocsex{daru}, de \cirilocsex{darU}) ``presente'' com, no
\emph{Saltério}, uma desinência mista de \gloss{nom.pl.} \emph{-ovi}, refação de
\emph{-ove} alinhado com o \emph{-i} dos temas em \emph{-o-}). O
\gloss{inst.sg.} -\emph{ъmь} do tipo em -\emph{u}- é encontrado também por vezes
no tipo em -\emph{o}-, o que pode refletir um moravismo, dado que
\emph{-ъmь} é regular no tipo em -\emph{o}- no eslavo ocidental e no
eslavo oriental, em contraste com o eslavo meridional que tem
-\emph{omь}. A contaminação entre temas em -\emph{o}- e temas em
-\emph{u}- é já do eslavo comum, e o \gloss{nom.sg.} -\emph{ъ} dos temas em
\emph{-o}- é, ele mesmo, um empréstimo ao tipo em -\emph{u}-.
\clearpage

Os temas em -\emph{u}- eram paralelos aos temas em -\emph{i-} em
indo-europeu. Esse paralelismo foi largamente apagado pelas evoluções
fonéticas próprias do eslavo, notadamente nos resultados dos ditongos.

% TeX root=../main.tex

\begin{table}
  \caption[Temas em \ocsex{-i-}]{Temas em \ocsex{-i-} (Paradigma: \gloss{fem.}~\ocsex{kostI} `boca', \gloss{masc.}~\ocsex{põtI} `caminho')}\label{tab:temasemi}
	\begin{center}
		\begin{tabular}[c]{lllllll}
			\toprule
			                               & \multicolumn{2}{c}{Singular} & \multicolumn{2}{c}{Dual} & \multicolumn{2}{c}{Plural}                                                                                    \\
			\midrule
			\emph{nom.}                    & \cirilex{kostI}              & \ocsex{kost-I}           & \cirilex{kosti}            & \ocsex{kost-i}   & \cirilex{kosti}                & \ocsex{kost-i}               \\
			                               &                              &                          &                            &                  & \cirilex{põtije}               & \ocsex{põt-ije}              \\
			\rowcolor{gray!10} \emph{acu.} & \cirilex{kostI}              & \ocsex{kost-I}           & \cirilex{kosti}            & \ocsex{kost-i}   & \cirilex{kosti}                & \ocsex{kost-i}               \\
			\emph{gen.}                    & \cirilex{kosti}              & \ocsex{kost-i}           & \cirilex{kostiju}          & \ocsex{kost-iju} & \cirilex{kostii}               & \ocsex{kost-ijI}             \\
			\rowcolor{gray!10} \emph{dat.} & \cirilex{kosti}              & \ocsex{kost-i}           & \cirilex{kostIma}          & \ocsex{kost-Ima} & \cirilex{kostImU}              & \ocsex{kost-ImU}             \\
			\rowcolor{gray!10}             &                              &                          &                            &                  & \hspace{0.5em}[\cirilex{-emU}] & \hspace{0.5em}[\ocsex{-emU}] \\
			\emph{inst.}                   & \cirilex{kostijõ}            & \ocsex{kost-ijõ}         & \cirilex{kostIma}          & \ocsex{kost-Ima} & \cirilex{kostImi}              & \ocsex{kost-Imi}             \\
			                               & \cirilex{põtImI}             & \ocsex{põt-ImI}          &                            &                  &                                &                              \\
			\rowcolor{gray!10} \emph{lac.} & \cirilex{kosti}              & \ocsex{kost-i}           & \cirilex{kostiju}          & \ocsex{kost-iju} & \cirilex{kostIxU}              & \ocsex{kost-IxU}             \\
			\rowcolor{gray!10}             &                              &                          &                            &                  & \hspace{0.5em}[\cirilex{-exU}] & \hspace{0.5em}[\ocsex{-exU}] \\
			\emph{voc.}                    & \cirilex{kosti}              & \ocsex{kost-i}           &                            &                  &                                &                              \\
			\bottomrule
		\end{tabular}
	\end{center}
\end{table}


O tipo em -\emph{i}- compreende sobretudo femininos, e é produtivo
apenas para formar abstratos femininos, o que continua uma das funções
do sufixo indo-europeu *-\emph{i}-. Os raros masculinos como \emph{pǫtь}
ou \emph{l'udije} ``a gente/as gentes'' (\emph{pluralia tantum}) estão
em vias de serem absorvidos pelo tipo em -\emph{jo}- mole (\emph{kon'ь}
``cavalo''). A passagem de um tipo a outro é facilitada pela identidade
de várias desinências nos dois tipos (\gloss{nom.acu.sg.} -\emph{ь}, \gloss{loc.sg.}
-\emph{i}).

No \gloss{inst.sg.} aparece uma diferenciação de gênero, dado que os masculinos
em -\emph{i}- têm a desinência etimológica -\emph{ьmь}, às vezes o
alomorfe -\emph{emь} do tipo em -\emph{jo}- (\emph{kon'emь}), enquanto
os femininos têm -\emph{ьjǫ} (notado -\cirilex{ijõ} ou -\cirilex{Ijõ}), de
origem pronominal,
como os temas em -\emph{a}- (\emph{ženojǫ}). O mesmo ocorre no \gloss{nom.pl.}
(\gloss{masc.} \emph{-ije}, \gloss{fem.} -\emph{i}), mas para os outros casos o
masculino e o feminino têm a mesma forma. Com exceção do \gloss{inst.sg.}, esse
tipo é caracterizado por uma forma única para os casos oblíquos do
singular, consequência da lei das sílabas abertas (\gloss{gen.sg.} -\emph{i}
< *-\emph{ei̯s}, \gloss{dat.sg.} -\emph{i} < *-\emph{ei̯ei̯} com
haplologia, \gloss{loc.sg.} -\emph{i} < *-\emph{ēi̯}); essa mesma forma
é, de resto, idêntica à dos casos retos do plural. Os alomorfes
-\emph{emъ} e -\emph{exъ} no \gloss{dat.pl.} e no \gloss{loc.pl.} são derivadas de temas
consonantais e são regulares nos textos de proveniência macedônia, os
textos de proveniência búlgara conservam melhor as desinências próprias
dos temas em -\emph{i­­}-, -\emph{ьmь}, -\emph{ьmъ} e \emph{-ьxъ}. A
flutuação entre as desinências próprias do tipo em -\emph{i}-, em
-\emph{ь}-, e as desinências em -\emph{e}- têm também uma causa
fonética: \emph{ь}, em posição forte diante de \emph{jer} final, se
pronunciava [e] em macedônio antigo e em uma parte do búlgaro antigo
(\emph{cf.}~\ref{sec:jers}), o que faz com que a distinção entre as duas séries de
desinências aparecesse como um fato gráfico antes de tudo; como a
confusão das duas séries se faz em prol de \emph{-emь}, -\emph{emъ},
-\emph{exъ}, é normal que essa série seja a mais frequente. O mesmo vale
para os temas consonantais, que apresentam os mesmos alomorfes para
essas desinências.



\subsubsection{Temas consonantais}

Os temas consonantais, em desuso, estão em vias de serem absorvidos pelo
tipo em -\emph{i}- e, por isso, tem duas desinências nas formas de
vários casos, uma etimológica e outra emprestada do tipo em -\emph{i}-.
A passagem dos temas consonantais aos temas em -\emph{i}- é observada
também em báltico; é uma consequência dos resultados das nasais
vocálicas (\gloss{acu.sg.} consonantal *-\emph{m̥} >
esl.~\emph{*-ĭm}, balt. *-\emph{im}, homófono do \gloss{acu.sg.} *-\emph{im} dos temas
em -\emph{i}-, e por isso há a refação do \gloss{nom.sg.} em *-\emph{is} e a
passagem para o tipo em -\emph{i}-).

Os temas consonantais contêm vários subtipos, definidos pela natureza da
consoante final do tema. Eles são caracterizados pela homofonia do
\gloss{gen.sg.} e \gloss{loc.sg;}, como é o caso também dos temas em -\emph{u}- e em
-\emph{i}-. A única diferença entre os subtipos é a forma do \gloss{nom.sg.} que
não é dedutível do tema da palavra; para o resto a flexão é idêntica. O
isolamento da forma do \gloss{nom.sg.} no interior do paradigma é consequência da
lei das sílabas abertas, que fez desaparecer a consoante final do tema
na ausência de uma desinência vocálica, seguida dos resultados
específicos em sílaba final para alguns tipos. Como os tipos
consonantais estavam em desuso, eles foram influenciados em diferentes
graus pelos tipos vocálicos em -\emph{o}-, -\emph{a}- e -\emph{i}-.

Os temas em nasal (\autoref{tab:temasemnasal}) são ainda bastante numerosos.
Esse tipo contém masculinos e neutros.
Os neutros em -\emph{mę} continuam o
*\emph{-mn̥} original indo-europeu; os masculinos do tipo \emph{kamy}
continuam o tipo em \emph{*-mōn}, *-\emph{men-}
(\glsxtrshort{lituano}~\emph{akmuō}, \gloss{gen.}
\emph{akmen̄s} ``pedra'', \glsxtrshort{grego}~ἄκμων ``bigorna'').
O tipo masculino em
-\emph{n}- tem duas variantes no \gloss{nom.sg.}, o almorfe -\emph{y} do tipo
\emph{kamy} (< *-\emph{nos} < *-\emph{ōn-s},
recaracterização de *-\emph{ōn} original, preferível à explicação que
parte diretamente da forma herdada *-\emph{ōn} e supõe um fechamento em
final *-\emph{ōn}\# > *-\emph{ūn} > ­-\emph{y}
contradita pelo tratamento da \gloss{1ªsg} *-\emph{ōm}\# >
*-\emph{ōn} > -\emph{ǫ}) e um alomorfe -\emph{ę}
(< -\emph{ēn(s)}, do tipo \glsxtrshort{grego}~ἄρρην ``macho'') no subtipo
\gloss{nom.sg.}
\cirilocsex{korẽ}, \gloss{acu.sg.} \cirilocsex{korenI} ``raiz''; quanto ao resto
da flexão, é idêntica nos dois subtipos. Nos dois casos, a forma própria
do \gloss{nom.sg.} sofre concorrência da forma em -\emph{enь} idêntica ao
\gloss{acu.sg.}.
\clearpage

% TeX root=../main.tex

\begin{table}[!htb]
	\begin{adjustwidth}{-2cm}{0cm}
		\caption[Temas em nasal]{Temas em nasal (Paradigmas:
			\gloss{masc.}~\ocsex{dInI} `dia', \ocsex{kamy} `pedra', \ocsex{jelenI} `cervo';
			\gloss{neut.}~\ocsex{imẽ} `nome')}\label{tab:temasemnasal}
		\begin{center}
			\begin{tabular}[c]{llllllll}
				\toprule
				                   &                  & \multicolumn{2}{c}{masculino radical} & \multicolumn{2}{c}{masculino sufixado} & \multicolumn{2}{c}{neutro}                                                            \\
				\midrule
				\emph{sg.}         & \emph{nom.}      & \cirilex{dInI}                        & \ocsex{dIn-I}                          & \cirilex{kamy}             & \ocsex{kamy}      & \cirilex{imẽ}     & \ocsex{imẽ}      \\
				                   &                  &                                       &                                        & \cirilex{kamenI}           & \ocsex{kamen-I}   &                   &                  \\
				\rowcolor{gray!10} & \emph{acu.}      & \cirilex{dInI}                        & \ocsex{dIn-I}                          & \cirilex{kamenI}           & \ocsex{kamen-I}   & \cirilex{imẽ}     & \ocsex{imẽ}      \\
				                   & \emph{gen.}      & \cirilex{dIne}                        & \ocsex{dIn-e}                          & \cirilex{kamene}           & \ocsex{kamen-e}   & \cirilex{imene}   & \ocsex{imen-e}   \\
				                   &                  & \cirilex{dIni}                        & \ocsex{dIn-i}                          & \cirilex{kameni}           & \ocsex{kamen-i}   & \cirilex{imeni}   & \ocsex{imen-i}   \\
				\rowcolor{gray!10} & \emph{dat.}      & \cirilex{dIni}                        & \ocsex{dIn-i}                          & \cirilex{kamenu}           & \ocsex{kamen-u}   & \cirilex{imenu}   & \ocsex{imen-u}   \\
				                   & \emph{inst.}     & \cirilex{dInImI}                      & \ocsex{dIn-ImI}                        & \cirilex{kamenemI}         & \ocsex{kamen-emI} & \cirilex{imenImI} & \ocsex{imen-ImI} \\
				                   &                  & [\cirilex{-emI}]                      & [\ocsex{-emI}]                         &                            &                   & [\cirilex{-emI}]  & [\ocsex{-emI}]   \\
				\rowcolor{gray!10} & \emph{loc.}      & \cirilex{dIne}                        & \ocsex{dIn-e}                          & \cirilex{kamene}           & \ocsex{kamen-e}   & \cirilex{imene}   & \ocsex{imen-e}   \\
				                   &                  & \cirilex{dIni}                        & \ocsex{dIn-i}                          & \cirilex{kameni}           & \ocsex{kamen-i}   & \cirilex{imeni}   & \ocsex{imen-i}   \\
				                   & \emph{voc.}      & \cirilex{dInI}                        & \ocsex{dIn-I}                          &                            &                   & \cirilex{imenu}   & \ocsex{imen-u}   \\
				\midrule
				\emph{du.}         & \emph{nom.acu}   & \cirilex{dIni}                        & \ocsex{dIn-i}                          & \cirilex{jeleni}           & \ocsex{jelen-i}   & \cirilex{imenE}   & \ocsex{imen-E}   \\
				                   &                  &                                       &                                        &                            &                   & \cirilex{imeni}   & \ocsex{imen-i}   \\
				\rowcolor{gray!10} & \emph{gen.loc.}  & \cirilex{dInu}                        & \ocsex{dIn-u}                          & \cirilex{jeleniju}         & \ocsex{jelen-iju} & \cirilex{imenu}   & \ocsex{imen-u}   \\
				\rowcolor{gray!10} &                  & \cirilex{dIniju}                      & \ocsex{dIn-iju}                        &                            &                   &                   &                  \\
				                   & \emph{dat.inst.} & \cirilex{dInIma}                      & \ocsex{dIn-Ima}                        & \cirilex{jeleIoma}         & \ocsex{jelen-Ima} & \cirilex{imenIma} & \ocsex{imen-Ima} \\
				\midrule
				\emph{pl.}         & \emph{nom.}      & \cirilex{dIne}                        & \ocsex{dIn-e}                          & \cirilex{jelene}           & \ocsex{jelen-e}   & \cirilex{imena}   & \ocsex{imen-a}   \\
				                   &                  & \cirilex{dInije}                      & \ocsex{dIn-ije}                        & \cirilex{jelenije}         & \ocsex{jelen-ije} &                   &                  \\
				\rowcolor{gray!10} & \emph{acu.}      & \cirilex{dIni}                        & \ocsex{dIn-i}                          & \cirilex{jeleni}           & \ocsex{jelen-i}   & \cirilex{imena}   & \ocsex{imen-a}   \\
				                   & \emph{gen.}      & \cirilex{dInU}                        & \ocsex{dIn-U}                          & \cirilex{jelenU}           & \ocsex{jelen-U}   & \cirilex{imenU}   & \ocsex{imen-U}   \\
				                   &                  & \cirilex{dInii}                       & \ocsex{dIn-ijU}                        & \cirilex{jelenii}          & \ocsex{jelen-ijU} &                   &                  \\
				\rowcolor{gray!10} & \emph{dat.}      & \cirilex{dInImU}                      & \ocsex{dIn-ImU}                        & \cirilex{jelenemU}         & \ocsex{jelen-emU} & \cirilex{imenemU} & \ocsex{imen-emU} \\
				\rowcolor{gray!10} &                  & [\cirilex{-emU}]                      & [\ocsex{-emU}]                         &                            &                   &                   &                  \\
				                   & \emph{inst.}     & \cirilex{dIny}                        & \ocsex{dIn-y}                          & \cirilex{jelenImi}         & \ocsex{jelen-Imi} & \cirilex{imeny}   & \ocsex{imen-y}   \\
				                   &                  & \cirilex{dInImi}                      & \ocsex{dIn-Imi}                        &                            &                   &                   &                  \\
				\rowcolor{gray!10} & \emph{loc.}      & \cirilex{dInIxU}                      & \ocsex{dIn-IxU}                        & \cirilex{jelenExU}         & \ocsex{jelen-ExU} & \cirilex{imenixU} & \ocsex{imen-ixU} \\
				\rowcolor{gray!10} &                  & [\cirilex{-exU}]                      & [\ocsex{-exU}]                         &                            &                   &                   &                  \\
				\bottomrule
			\end{tabular}
		\end{center}
	\end{adjustwidth}
\end{table}


Os \emph{pluralia tantum} em -\emph{{(j)}ane} (masculinos) se
declinam parcialmente segundo o modelo (\gloss{nom.} \cirilocsex{graZdane}
``cidadãos'', \gloss{acu.} -\emph{any}, \gloss{gen.} \emph{-anъ}, \gloss{dat.}
\emph{-anemъ}, \gloss{inst.} \emph{-any}, \gloss{loc.} -\emph{anexъ}).
Seu singular é fornecido por um singulativo em -\emph{inъ}
(\cirilocsex{graZdaninU}) que segue o tipo
em -\emph{o}- (\cirilocsex{gradU}).

Os outros temas consonantais em ressoante (\autoref{tab:temasemliquidaeuu}) são
mais raros e contêm apenas poucas unidades cada um.
\clearpage

% TeX root=../main.tex

\begin{table}[!htb]
	\begin{adjustwidth}{0cm}{-2.2cm}
		\caption[Temas em líquida e -\emph{ū}-]{Temas em líquida e -\emph{ū}- (Paradigmas:
			\gloss{fem.}~\ocsex{mati} `mãe', \ocsex{crIky} `igreja';
			\gloss{masc.}~\ocsex{dElatel'e}
			`trabalhadores')}\label{tab:temasemliquidaeuu}
		\begin{center}
			\begin{tabular}[c]{llllllll}
				\toprule
				                   &              & \multicolumn{2}{c}{femininos em \emph{-r-}} & \multicolumn{2}{c}{masculinos em -\emph{l}- e -\emph{r}-} & \multicolumn{2}{c}{femininos em -\emph{ū}-}                                                                  \\
				\midrule
				\emph{sg.}         & \emph{nom.}  & \cirilex{mati}                              & \ocsex{mati}                                              &                                             &                     & \cirilex{crIky}     & \ocsex{crIky}      \\
				                   &              &                                             &                                                           &                                             &                     &                     &                    \\
				\rowcolor{gray!10} & \emph{acu.}  & \cirilex{materI}                            & \ocsex{mater-I}                                           &                                             &                     & \cirilex{crIkUvI}   & \ocsex{crIkUv-I}   \\
				                   & \emph{gen.}  & \cirilex{matere}                            & \ocsex{mater-e}                                           &                                             &                     & \cirilex{crIkUve}   & \ocsex{crIkUv-e}   \\
				                   &              & \cirilex{materi}                            & \ocsex{mater-i}                                           &                                             &                     & \cirilex{crIkUvi}   & \ocsex{crIkUv-i}   \\
				\rowcolor{gray!10} & \emph{dat.}  & \cirilex{materi}                            & \ocsex{mater-i}                                           &                                             &                     & \cirilex{crIkUvijõ} & \ocsex{crIkUv-ijõ} \\
				                   & \emph{inst.} & \cirilex{materijõ}                          & \ocsex{mater-ijõ}                                         &                                             &                     & \cirilex{crIkUvImI} & \ocsex{crIkUv-ImI} \\
				\rowcolor{gray!10} & \emph{loc.}  & \cirilex{materi}                            & \ocsex{mater-i}                                           &                                             &                     & \cirilex{crIkUve}   & \ocsex{crIkUv-e}   \\
				\rowcolor{gray!10} &              &                                             &                                                           &                                             &                     & \cirilex{crIkUvi}   & \ocsex{crIkUv-i}   \\
				                   & \emph{voc.}  & \cirilex{mati}                              & \ocsex{mati}                                              &                                             &                     & \cirilex{crIky}     & \ocsex{crIky}      \\
				\midrule
				\emph{pl.}         & \emph{nom.}  & \cirilex{materi}                            & \ocsex{mater-i}                                           & \cirilex{dElatel'e}                         & \ocsex{dElatel'e}   & \cirilex{crIkUvi}   & \ocsex{crIkUv-i}   \\
				\rowcolor{gray!10} & \emph{acu.}  & \cirilex{materi}                            & \ocsex{mater-i}                                           & \cirilex{dElatel'i}                         & \ocsex{dElatel'i}   & \cirilex{crIkUvi}   & \ocsex{crIkUv-i}   \\
				                   & \emph{gen.}  & \cirilex{materU}                            & \ocsex{mater-U}                                           & \cirilex{dElatel'U}                         & \ocsex{dElatel'U}   & \cirilex{crIkUvU}   & \ocsex{crIkUv-U}   \\
				                   &              & \cirilex{materii}                           & \ocsex{mater-ijU}                                         & \cirilex{dElatel'ii}                        & \ocsex{dElatel'ijU} &                     &                    \\
				\rowcolor{gray!10} & \emph{dat.}  & \cirilex{materemU}                          & \ocsex{mater-emU}                                         & \cirilex{dElatel'emU}                       & \ocsex{dElatel'emU} & \cirilex{crIkUvamU} & \ocsex{crIkUv-amU} \\
				\rowcolor{gray!10} &              &                                             &                                                           & \cirilex{dElatel'imU}                       & \ocsex{dElatel'imU} &                     &                    \\
				                   & \emph{inst.} & \cirilex{materImi}                          & \ocsex{mater-Imi}                                         & \cirilex{dElatel'y}                         & \ocsex{dElatel'y}   & \cirilex{crIkUvami} & \ocsex{crIkUv-ami} \\
				                   &              &                                             &                                                           & \cirilex{dElatel'i}                         & \ocsex{dElatel'i}   &                     &                    \\
				\rowcolor{gray!10} & \emph{loc.}  & \cirilex{materexU}                          & \ocsex{mater-exU}                                         & \cirilex{dElatel'ExU}                       & \ocsex{dElatel'ExU} & \cirilex{crIkUvaxU} & \ocsex{crIkUv-axU} \\
				\rowcolor{gray!10} &              &                                             &                                                           & \cirilex{dElatel'ixU}                       & \ocsex{dElatel'ExU} &                     &                    \\
				\bottomrule
			\end{tabular}
		\end{center}
	\end{adjustwidth}
\end{table}


O dual não é atestado para esses tipos flexionais.

Os femininos em -\emph{r}- contêm apenas dois substantivos \cirilocsex{mati}
``mãe'' (tema \emph{mater-}) e \cirilocsex{dUSti} ``filha'' (tema \emph{dъšter-}). 
Notar-se-á que o \gloss{acu.sg.} tem uma variante \cirilocsex{matere}, com a
forma do genitivo, ou seja, o genitivo-acusativo animado conseguiu espaço também
nesse tipo flexional (i.e., estendendo-se para além do masculino).

Os temas em -\emph{l}- conservam a flexão consonantal apenas no plural,
mas essa categoria é bem representada dada a frequência do sufixo de
nome de agente -\emph{tel'ь}. O singular dos masculinos em -\emph{tel'e}
segue o tipo em -\emph{jo}- (\emph{kon'ь}) e seu plural é igualmente
tomado do tipo em \emph{-jo}- (\gloss{acu.pl.} -\emph{ę}, \gloss{gen.pl.} -\emph{ь},
\gloss{loc.pl.} \emph{-ixъ}): as desinências do tipo consonantal são
minoritárias. Existe um tipo paralelo em -\emph{r'}- (\cirilocsex{mytar'e}
``publicanos''), cujo singular segue também o tipo em -\emph{jo}-.

O tipo em -\emph{y} (< *-\emph{ū} < *-\emph{uh\textsubscript{2}})
hesita no singular entre a flexão
consonantal e o tipo em -\emph{i}-; no plural ele já passou largamente
ao tipo em -\emph{a}- (\gloss{dat.inst.loc.pl.}). Um genitivo-acusativo é
igualmente atestado nesse tipo para os substantivos que designam pessoas
(\gloss{acu.sg.} \emph{svekrъve svojo} ``sua sogra'' \passageme{Mt1035} \glsxtrshort{Zogr}, em
comparação com \emph{svekrъvь svojǫ} em \glsxtrshort{Marianus}): o desenvolvimento
é paralelo à extensão do gen.acu. animado para \emph{mati} ``mãe''.


% TeX root=../main.tex

\begin{table}[!htb]
	\begin{adjustwidth}{0cm}{0cm}
		\caption[Temas neutros em oclusiva e em \emph{-s-}]{Temas neutros em oclusiva e em \emph{-s-} (Paradigmas:
			\ocsex{agnẽ} `cordeiro', \ocsex{slovo} `palavra')}\label{tab:temasemoclusivaes}
		\begin{center}
			\begin{tabular}[c]{llllll}
				\toprule
				                   &                  & \multicolumn{2}{c}{neutros em \emph{-t-}} & \multicolumn{2}{c}{neutros em \emph{-s-}}                                            \\
				\midrule
				\emph{sg.}         & \emph{nom.acu}   & \cirilex{agnẽ}                            & \ocsex{agnẽ}                              & \cirilex{slovo}     & \ocsex{slovo}      \\
				                   &                  &                                           &                                           & \cirilex{slovesI}   & \ocsex{sloves-I}   \\
				\rowcolor{gray!10} & \emph{gen.}      & \cirilex{agnẽte}                          & \ocsex{agnẽt-e}                           & \cirilex{slovese}   & \ocsex{sloves-e}   \\
				\rowcolor{gray!10} &                  &                                           &                                           & \cirilex{slovesi}   & \ocsex{sloves-i}   \\
				                   & \emph{dat.}      & \cirilex{agnẽti}                          & \ocsex{agnẽt-i}                           & \cirilex{slovesu}   & \ocsex{sloves-u}   \\
				                   & \emph{inst.}     & \cirilex{agnẽtemI}                        & \ocsex{agnẽt-emI}                         & \cirilex{slovesemI} & \ocsex{sloves-emI} \\
				                   &                  &                                           &                                           & \cirilex{slovesImI} & \ocsex{sloves-ImI} \\
				\rowcolor{gray!10} & \emph{loc.}      & \cirilex{agnẽte}                          & \ocsex{agnẽt-e}                           & \cirilex{slovese}   & \ocsex{sloves-e}   \\
				\rowcolor{gray!10} &                  & \cirilex{agnẽti}                          & \ocsex{agnẽt-i}                           & \cirilex{slovesi}   & \ocsex{sloves-i}   \\
				\midrule
				\emph{du.}         & \emph{nom.acu}   & \cirilex{agnẽtE}                          & \ocsex{agnẽt-E}                           & \cirilex{slovesE}   & \ocsex{sloves-E}   \\
				                   &                  &                                           &                                           & \cirilex{slovesi}   & \ocsex{sloves-i}   \\
				                   &                  &                                           &                                           &                     &                    \\
				\rowcolor{gray!10} & \emph{gen.loc.}  & \cirilex{agnẽtu}                          & \ocsex{agnẽt-u}                           & \cirilex{slovesu}   & \ocsex{sloves-u}   \\
				                   & \emph{dat.inst.} & \cirilex{agnẽtIma}                        & \ocsex{agnẽt-Ima}                         & \cirilex{slovosIma} & \ocsex{sloves-Ima} \\
				\midrule
				\emph{pl.}         & \emph{nom.acu}   & \cirilex{agnẽta}                          & \ocsex{agnẽt-a}                           & \cirilex{slovesa}   & \ocsex{sloves-a}   \\
				\rowcolor{gray!10} & \emph{gen.}      & \cirilex{agnẽtU}                          & \ocsex{agnẽt-U}                           & \cirilex{slovesU}   & \ocsex{sloves-U}   \\
				                   & \emph{dat.}      & \cirilex{agnẽtemU}                        & \ocsex{agnẽt-emU}                         & \cirilex{slovesemU} & \ocsex{sloves-emU} \\
				\rowcolor{gray!10} & \emph{inst.}     & \cirilex{agnẽty}                          & \ocsex{agnẽt-y}                           & \cirilex{slovesy}   & \ocsex{sloves-y}   \\
				                   & \emph{loc.}      & \cirilex{agnẽtexU}                        & \ocsex{agnẽt-exU}                         & \cirilex{slovesExU} & \ocsex{sloves-ExU} \\
				\bottomrule
			\end{tabular}
		\end{center}
	\end{adjustwidth}
\end{table}


O tipo em \emph{-ęt}- < *-\emph{ent}- (\autoref{tab:temasemoclusivaes}) contém
apenas uma dezena de termos, essencialmente nomes de pequenos animais.

O tipo neutro em -\emph{s}- (\autoref{tab:temasemoclusivaes}), que continua o tipo
indo-europeu do lat. \emph{genus}, \emph{generis}, gr. γένος, γένεος/γένους, do qual conserva a apofonia (o sufixo tem vocalismo \emph{o} nos
casos retos no singular e \emph{e} em todos os outros), foi tomado do
tipo neutro em -\emph{o}- (\cirilocsex{mEsto} ``lugar''), dado que os
dois paradigmas têm em comum a forma de \gloss{nom.acu.sg.} em -\emph{o}. Por
isso, encontram-se regularmente formas do tipo \gloss{gen.sg.}
\cirilocsex{slova} em vez de \cirilocsex{slovese}, \gloss{loc.sg.} \cirilocsex{tElE}
``no corpo'' em vez de \cirilocsex{tElese}, e isso ocorre tanto na
tradição búlgara como na macedônia. 
Inversamente, a flexão em -\emph{s}-
é encontrada por vezes para os neutros em -\emph{o}- etimológicos
(sobretudo no \glsxtrshort{Supr}: \gloss{gen.sg.} \cirilocsex{drEvese} ``da árvore'',
para \cirilocsex{drEva}, \gloss{nom.pl.} \cirilocsex{dElesa} ``os atos'' para
\cirilocsex{dEla}).

Pertencem a esse tipo dois substantivos irregulares, \cirilocsex{oko}
``olho'' e \cirilocsex{uxo} ``orelha'', que seguem o tipo em -\emph{s}-
no singular e no plural (esse último é raríssimo), por apresentarem uma
palatalização (\gloss{gen.sg.} \cirilocsex{oCese}, \cirilocsex{uSese}) e cujo
dual, usual já que trata-se de pares naturais, é uma forma radical não
um tema em -\emph{s}-: \gloss{nom.acu.du.} \cirilocsex{oCi}, \cirilocsex{uSi} (que
continuam diretamente os duais indo-europeus
*\pietrans{h3ekw-ih1}
 e
*\pietrans{h2ews-ih1}), \gloss{gen.loc.}
\cirilocsex{oCiju}, \cirilocsex{uSiju}, \gloss{dat.inst.} \cirilocsex{oCima},
\cirilocsex{uSima}.

\section{Adjetivos e particípios}\label{sec:adjetivos}
\subsection{Paradigmas}\label{sec:adjparadigma}

Os adjetivos são marcados para gênero, número e caso: essas categorias
não são inerentes, o adjetivo simplesmente reproduz as categorias do
substantivo com que se relaciona e com o qual concorda. Um adjetivo
substantivado reproduz as categorias de seu referente implícito. O mesmo
ocorre com os particípios, formas adjetivais do verbo.

Os adjetivos têm flexão dupla em AEE, com uma forma determinada, que é
empregue quando o adjetivo qualifica um substantivo determinado ou
quando ele próprio é substantivado (\emph{cf.}~\ref{sec:determinação}), e uma forma
indeterminada, que é empregue nos outros contextos. A forma determinada
é o termo marcado desse par. A flexão determinada do adjetivo é uma
inovação comum ao eslavo e ao báltico.

A flexão indeterminada (\autoref{tab:adjetivosindeterminados}), chamada ``forma curta'' ou ``forma
simples'' segue a flexão dos substantivos: para o tipo duro, \gloss{masc.}~\cirilocsex{novU}, \gloss{fem.} \cirilocsex{nova}, \gloss{nt.} \cirilocsex{novo}
``novo'',
seguem respectivamente a flexão de \cirilocsex{gradU}, \cirilocsex{Zena},
\cirilocsex{mEsto}; para o tipo mole, \gloss{masc.}~\cirilocsex{niStU},
\gloss{fem.}~\cirilocsex{niSta}, \gloss{nt.}~\cirilocsex{niSte} seguem a
flexão de \cirilocsex{kon'I}, \cirilocsex{duSa}, \cirilocsex{niSta},
\cirilocsex{mor'e}. Os temas em velar apresentam a mesma alternância que os
substantivos: \emph{k $\sim$ c, g $\sim$ dz, x
$\sim$ s} (\cirilocsex{velikU} ``grande'',
\gloss{loc.sg.masc.}~\cirilocsex{velicE}, \cirilocsex{mUnogU} ``abundante'',
\gloss{nom.pl.masc.}~\cirilocsex{mUnodz'i,mUnodzi} ``numerosos'').

% TeX root=../main.tex
\begin{table}[!htb]
	\begin{adjustwidth}{0cm}{0cm}
		\caption[Flexão indeterminada]{Flexão indeterminada (Paradigmas:
			\ocsex{novU} `novo', \ocsex{niStI}
			`pobre')}\label{tab:adjetivosindeterminados}
		\begin{center}
			\begin{tabular}[c]{llllllll}
				\toprule
				                   &                   & \multicolumn{2}{c}{masc.neut. duro} & \multicolumn{2}{c}{fem. duro} & \multicolumn{2}{c}{masc.neut. mole}                                                            \\
				\midrule
				\emph{sg.}         & \emph{nom.}       & \cirilex{novU}                      & \ocsex{nov-U}                 & \cirilex{nova}                      & \ocsex{nov-a}   & \cirilex{niStI}    & \ocsex{niSt-I}    \\
				\rowcolor{gray!10} & \emph{acu.}       & \cirilex{novU/-a}                   & \ocsex{nov-U/-a}              & \cirilex{novõ}                      & \ocsex{nov-õ}   & \cirilex{niStI/-a} & \ocsex{niSt-I/-a} \\
				                   & \emph{n.-a.neut.} & \cirilex{novo}                      & \ocsex{nov-o}                 &                                     &                 & \cirilex{niSte}    & \ocsex{niSt-e}    \\
				\rowcolor{gray!10} & \emph{gen.}       & \cirilex{nova}                      & \ocsex{nov-a}                 & \cirilex{novy}                      & \ocsex{nov-y}   & \cirilex{niSta}    & \ocsex{niSt-a}    \\
				                   & \emph{dat.}       & \cirilex{novu}                      & \ocsex{nov-u}                 & \cirilex{novE}                      & \ocsex{nov-E}   & \cirilex{niStu}    & \ocsex{niSt-u}    \\
				\rowcolor{gray!10} & \emph{inst.}      & \cirilex{novomI}                    & \ocsex{nov-omI}               & \cirilex{novojõ}                    & \ocsex{nov-ojõ} & \cirilex{niStemI}  & \ocsex{niSt-emI}  \\
				                   & \emph{loc.}       & \cirilex{novE}                      & \ocsex{nov-E}                 & \cirilex{novE}                      & \ocsex{nov-E}   & \cirilex{niSti}    & \ocsex{niSt-i}    \\
				\rowcolor{gray!10} & \emph{voc.}       & \cirilex{nove}                      & \ocsex{nov-e}                 & \cirilex{novo}                      & \ocsex{nov-o}   & \cirilex{niStu}    & \ocsex{niSt-u}    \\
				\midrule
				\emph{du.}         & \emph{nom.acu}    & \cirilex{nova}                      & \ocsex{nov-a}                 & \cirilex{novE}                      & \ocsex{nov-E}   & \cirilex{niSta}    & \ocsex{niSt-a}    \\
				\rowcolor{gray!10} & \emph{n.-a.neut.} & \cirilex{novE}                      & \ocsex{nov-E}                 &                                     &                 & \cirilex{niSti}    & \ocsex{niSt-i}    \\
				                   & \emph{gen.loc.}   & \cirilex{novu}                      & \ocsex{nov-u}                 & \cirilex{novu}                      & \ocsex{nov-u}   & \cirilex{niStu}    & \ocsex{niSt-u}    \\
				\rowcolor{gray!10} & \emph{dat.inst.}  & \cirilex{novoma}                    & \ocsex{nov-oma}               & \cirilex{novama}                    & \ocsex{nov-ama} & \cirilex{niStema}  & \ocsex{niSt-ema}  \\
				\midrule
				\emph{pl.}         & \emph{nom.}       & \cirilex{novi}                      & \ocsex{nov-i}                 & \cirilex{novy}                      & \ocsex{nov-y}   & \cirilex{niSti}    & \ocsex{niSt-i}    \\
				\rowcolor{gray!10} & \emph{acu.}       & \cirilex{novy}                      & \ocsex{nov-y}                 & \cirilex{novy}                      & \ocsex{nov-y}   & \cirilex{niStẽ}    & \ocsex{niSt-ẽ}    \\
				                   & \emph{n.-a.neut.} & \cirilex{nova}                      & \ocsex{nov-a}                 &                                     &                 & \cirilex{niSta}    & \ocsex{niSt-a}    \\
				\rowcolor{gray!10} & \emph{gen.}       & \cirilex{novU}                      & \ocsex{nov-U}                 & \cirilex{novU}                      & \ocsex{nov-U}   & \cirilex{niStI}    & \ocsex{niSt-I}    \\
				                   & \emph{dat.}       & \cirilex{novomU}                    & \ocsex{nov-omU}               & \cirilex{novamU}                    & \ocsex{nov-amU} & \cirilex{niStemU}  & \ocsex{niSt-emU}  \\
				\rowcolor{gray!10} & \emph{inst.}      & \cirilex{novy}                      & \ocsex{nov-y}                 & \cirilex{novami}                    & \ocsex{nov-ami} & \cirilex{niSti}    & \ocsex{niSt-i}    \\
				                   & \emph{loc.}       & \cirilex{novExU}                    & \ocsex{nov-ExU}               & \cirilex{novaxU}                    & \ocsex{nov-axU} & \cirilex{niStixU}  & \ocsex{niSt-ixU}  \\
				\bottomrule
			\end{tabular}
		\end{center}
	\end{adjustwidth}
\end{table}


A flexão determinada (\autoref{tab:adjetivosdeterminados}), chamada ``forma
longa'' ou ``forma
composta'', é caracterizada por desinências específicas e segue a flexão
pronominal; de fato, originalmente, a flexão determinada deriva do
ajuste de um tema pronominal enclítico à flexão indeterminada. Esse tema
é, etimologicamente, o do pronome relativo (\emph{cf.}~\ref{sec:anafóricoerelativo}).
Consequentemente, a flexão determinada tem uma dupla marca flexional, já
que o pronome enclítico flexionado se ajusta à forma do adjetivo já
flexionado. Contudo, a evolução fonética e, sobretudo, as contrações
advindas da queda do /j/ intervocálico fizeram com que várias
desinências já não fossem analisáveis e se tornassem morfemas complexos
específicos.

% TeX root=../main.tex
\begin{table}[!htb]
	\begin{adjustwidth}{0cm}{0cm}
		\caption[Flexão determinada]{Flexão determinada (Paradigmas: \ocsex{novU} `novo', \ocsex{niStI} `pobre')}\label{tab:adjetivosdeterminados}
		\begin{center}
			\begin{tabular}[c]{llllllll}
				\toprule
				                   &                   & \multicolumn{2}{c}{masc.neut. duro}    & \multicolumn{2}{c}{fem. duro}         & \multicolumn{2}{c}{masc.neut. mole}                                                                                              \\
				\midrule
				\emph{sg.}         & \emph{nom.}       & \cirilex{novyI}                        & \ocsex{nov-yjI}                       & \cirilex{novaja}                    & \ocsex{nov-aja}    & \cirilex{niStI}                    & \ocsex{niSt-I}                   \\
				\rowcolor{gray!10} & \emph{acu.}       & \cirilex{novyi}                        & \ocsex{nov-yjI}                       & \cirilex{novõjõ}                    & \ocsex{nov-õjõ}    & \cirilex{niStyI}                   & \ocsex{niSt-yjI}                 \\
				\rowcolor{gray!10} &                   & \hspace{1em}\slash\cirilex{-ajego}     & \hspace{1em}\slash\ocsex{-ajego}      &                                     &                    & \hspace{1em}\slash\cirilex{-ajego} & \hspace{1em}\slash\ocsex{-ajego} \\
				                   & \emph{n.-a.neut.} & \cirilex{novoje}                       & \ocsex{nov-oje}                       &                                     &                    & \cirilex{niSteje}                  & \ocsex{niSt-eje}                 \\
				\rowcolor{gray!10} & \emph{gen.}       & \cirilex{novajego}                     & \ocsex{nov-ajego}                     & \cirilex{novyjẽ}                    & \ocsex{nov-yjẽ}    & \cirilex{niStajego}                & \ocsex{niSt-ajego}               \\
				\rowcolor{gray!10} &                   & \cirilex{nova(a)go}                    & \ocsex{nov-a(a)go}                    &                                     &                    & \cirilex{niSta(a)go}               & \ocsex{niSt-a(a)go}              \\
				                   & \emph{dat.}       & \cirilex{novujemu}                     & \ocsex{nov-ujemu}                     & \cirilex{novEji}                    & \ocsex{nov-Eji}    & \cirilex{niStujemu}                & \ocsex{niSt-ujemu}               \\
				                   &                   & \cirilex{novu(u)mu}                    & \ocsex{nov-u(u)mu}                    &                                     &                    & \cirilex{niStu(u)mu}               & \ocsex{niSt-u(u)mu}              \\
				\rowcolor{gray!10} & \emph{inst.}      & \cirilex{novyimI}                      & \ocsex{nov-yjimI}                     & \cirilex{novõjõ}                    & \ocsex{nov-õjõ}    & \cirilex{niStiimI}                 & \ocsex{niSt-ijimI}               \\
				\rowcolor{gray!10} &                   & \cirilex{novymI}                       & \ocsex{nov-ymI}                       &                                     &                    & \cirilex{niStimI}                  & \ocsex{niSt-imI}                 \\
				                   & \emph{loc.}       & \cirilex{novEjemI}                     & \ocsex{nov-EjemI}                     & \cirilex{novEji}                    & \ocsex{nov-Eji}    & \cirilex{niStiimI}                 & \ocsex{niSt-ijimI}               \\
				                   &                   & \cirilex{novEmI}                       & \ocsex{nov-EmI}                       &                                     &                    & \cirilex{niStimI}                  & \ocsex{niSt-imI}                 \\
				\midrule
				\emph{du.}         & \emph{nom.acu}    & \cirilex{novaja}                       & \ocsex{nov-aja}                       & \cirilex{novEji}                    & \ocsex{nov-Eji}    & \cirilex{niStaja}                  & \ocsex{niSt-aja}                 \\
				\rowcolor{gray!10} & \emph{n.-a.neut.} & \cirilex{novEi}                        & \ocsex{nov-Eji}                       &                                     &                    & \cirilex{niStii}                   & \ocsex{niSt-iji}                 \\
				                   & \emph{gen.loc.}   & \multicolumn{2}{c}{\cirilex{novuju}}   & \multicolumn{2}{c}{\ocsex{nov-uju}}   & \cirilex{niStuju}                   & \ocsex{niSt-uju}                                                                           \\
				\rowcolor{gray!10} & \emph{dat.inst.}  & \multicolumn{2}{c}{\cirilex{novyjima}} & \multicolumn{2}{c}{\ocsex{nov-yjima}} & \cirilex{niStijima}                 & \ocsex{niSt-ijima}                                                                         \\
				\rowcolor{gray!10} &                   & \multicolumn{2}{c}{\cirilex{novyma}}   & \multicolumn{2}{c}{\ocsex{nov-yma}}   & \cirilex{niStima}                   & \ocsex{niSt-ima}                                                                           \\
				\midrule
				\emph{pl.}         & \emph{nom.}       & \cirilex{noviji}                       & \ocsex{nov-iji}                       & \cirilex{novyjẽ}                    & \ocsex{nov-yjẽ}    & \cirilex{niStiji}                  & \ocsex{niSt-iji}                 \\
				\rowcolor{gray!10} & \emph{acu.}       & \cirilex{novyjẽ}                       & \ocsex{nov-yjẽ}                       & \cirilex{novyjẽ}                    & \ocsex{nov-yjẽ}    & \cirilex{niStẽjẽ}                  & \ocsex{niSt-ẽjẽ}                 \\
				                   & \emph{n.-a.neut.} & \cirilex{novaja}                       & \ocsex{nov-aja}                       &                                     &                    & \cirilex{niStaja}                  & \ocsex{niSt-aja}                 \\
				\rowcolor{gray!10} & \emph{gen.}       & \multicolumn{2}{c}{\cirilex{novyjixU}} & \multicolumn{2}{c}{\ocsex{nov-yjixU}} & \cirilex{niStijixU}                 & \ocsex{niSt-ijixU}                                                                         \\
				\rowcolor{gray!10} &                   & \multicolumn{2}{c}{\cirilex{novyxU}}   & \multicolumn{2}{c}{\ocsex{nov-yxU}}   & \cirilex{niStixU}                   & \ocsex{niSt-ixU}                                                                           \\
				                   & \emph{dat.}       & \multicolumn{2}{c}{\cirilex{novyjimU}} & \multicolumn{2}{c}{\ocsex{nov-yjimU}} & \cirilex{niStijimU}                 & \ocsex{niSt-ijimU}                                                                         \\
				                   &                   & \multicolumn{2}{c}{\cirilex{novymU}}   & \multicolumn{2}{c}{\ocsex{nov-ymU}}   & \cirilex{niStimU}                   & \ocsex{niSt-imU}                                                                           \\
				\rowcolor{gray!10} & \emph{inst.}      & \multicolumn{2}{c}{\cirilex{novyjimi}} & \multicolumn{2}{c}{\ocsex{nov-yjimi}} & \cirilex{niStijimi}                 & \ocsex{niSt-ijimi}                                                                         \\
				\rowcolor{gray!10} &                   & \multicolumn{2}{c}{\cirilex{novymi}}   & \multicolumn{2}{c}{\ocsex{nov-ymi}}   & \cirilex{niStimi}                   & \ocsex{niSt-imi}                                                                           \\
				                   & \emph{loc.}       & \multicolumn{2}{c}{\cirilex{novyjixU}} & \multicolumn{2}{c}{\ocsex{nov-yjixU}} & \cirilex{niStijixU}                 & \ocsex{niSt-ijixU}                                                                         \\
				                   &                   & \multicolumn{2}{c}{\cirilex{novyxU}}   & \multicolumn{2}{c}{\ocsex{nov-yxU}}   & \cirilex{niStixU}                   & \ocsex{niSt-ixU}                                                                           \\
				\bottomrule
			\end{tabular}
		\end{center}
	\end{adjustwidth}
\end{table}


O tipo feminino mole é, por sua vez, baseado no tipo em -\emph{ja}-
(\emph{duša}) e não se distingue do tipo duro quanto às seguintes
formas: \gloss{gen.sg.} \cirilocsex{niStẽjẽ}, \gloss{dat.loc.sg.} \cirilocsex{niStiji},
\gloss{nom.acu.dual} \cirilocsex{niStiji}, \gloss{nom.acu.pl.} \cirilocsex{niStẽjẽ}.

A flexão determinada não possui forma própria de vocativo, no que difere
da flexão indeterminada (que tem o mesmo vocativo que os substantivos:
\emph{masc.} \cirilocsex{nove}, \gloss{fem.} \cirilocsex{novo}), é a forma de
nominativo que é empregue na função de vocativo. É característica da
flexão determinada a unificação de paradigmas do plural: os casos
oblíquos (\gloss{gen.}, \gloss{dat.}, \gloss{inst.}, \gloss{loc.}) são idênticos
para os três gêneros. 
A oposição de gênero é neutralizada em quase todo o paradigma do plural,
embora ela ainda exista para o plural da flexão indeterminada. Essa
neutralização é regular em todos os paradigmas pronominais e, de fato, a
flexão determinada do adjetivo é um tipo pronominal. O feminino
apresenta uma homofonia entre \gloss{acu.sg.} e \gloss{inst.sg.} no tipo duro e no mole
(desinência \cirilocsex{-õjõ} nos dois casos). O masculino apresenta,
por sua vez, uma homofonia entre \gloss{inst.sg.} e \gloss{loc.sg.} no tipo mole e os
dois casos continuam distintos no tipo duro. A aparente confusão no tipo
mole entre o \gloss{nom.sg.masc.} \cirilocsex{niStjI} e o \gloss{nom.pl.masc.}
\cirilocsex{niStiji} é apenas gráfica, as formas são distintas pela oposição
entre -\emph{ь} (\gloss{sg}) e -\emph{i} (\gloss{pl}) (cf.~\ref{sec:jers}).

Algumas formas casuais são ainda transparentes e revelam a origem da
forma determinada com sua dupla marca flexional: por exemplo,
\gloss{gen.sg.masc.} \cirilocsex{nova-jego}, acu.pl.masc. \cirilocsex{novy-jẽ},
\gloss{nom.sg.fem.} \cirilocsex{nova-ja}, \gloss{acu.sg.fem.}
\cirilocsex{novõ-jõ} etc.
Há, contudo, várias desinências que não são (mais) analisáveis seguindo
esse princípio.
Quase todo o plural foi formado por um tema secundário
\emph{novy-}, que ainda aparece no \gloss{inst.sg.masc.}, o qual não tem relação
alguma com sua forma indeterminada correspondente.
As formas contratas são majoritárias, sobretudo para o \gloss{gen.sg.masc.}
-\emph{aago/ -ago} e o
\gloss{dat.sg.masc.} \emph{-uumu} / \emph{-umu}, enquanto as formas não
contratas -\emph{ajego}, -\emph{ujemu} são raras. 
No \glsxtrshort{Sav} encontram-se já alguns exemplos de alinhamento na
desinência pronominal,
tendência que aparece nas línguas eslavas modernas: \gloss{gen.sg.masc.}
-\emph{ogo} (em vez de -\emph{ago}), \gloss{dat.sg.masc.} -\emph{omu} (em vez de
-\emph{umu}), conforme o demonstrativo \cirilocsex{togo}, \cirilocsex{tomu}.
No tipo mole, o \gloss{loc.pl.} determinado em \cirilocsex{-ixU}
advindo da contração de \cirilocsex{-ijixU} se confunde com o
\gloss{loc.pl.masc.nt.}, indeterminado (\cirilocsex{niStixU}, tipo
\cirilocsex{kon'ixU}, que, portanto, neutraliza a oposição entre forma
determinada e forma indeterminada, o que aparece também na flexão do
comparativo, do particípio presente ativo e do particípio passado ativo,
que correspondem ao tipo \cirilocsex{niStjU}. O \gloss{nom.sg.masc.} pode
apresentar uma variante contrata \cirilocsex{novy}.

Os temas em velar conservam na forma determinada a mesma alternância que
na forma indeterminada (\cirilocsex{velikyjI} ``grande'',
\gloss{loc.sg.masc.nt.} \cirilocsex{velicEjemI}, \cirilocsex{mUnogyjI}
``abundante'', \gloss{nom.pl.masc.} \cirilocsex{mUnodz'iji,mUnodziji}
``numerosos'').
\clearpage

\subsection{Graus de comparação}\label{sec:comparativo}

O adjetivo possui dois graus de comparação, um comparativo e um
superlativo. O comparativo (\autoref{tab:comparativo}) é de tipo sintético,
formado com
um sufixo específico -\emph{ьš-} (< *-\emph{is}-, grau Ø de
*-\emph{i̯os}-, cf. o superlativo grego -ισ-τος) que aparece em toda a
flexão exceto no \gloss{nom.acu.sg.} \gloss{masc.} e \gloss{nt.}.
Apenas \gloss{nom.sg.} e \gloss{acu.sg.}
masculino e neutro não têm esse sufixo sincronicamente (em diacronia, o
sufixo está presente mas a consoante final se emudece em virtude da lei
das sílabas abertas).


% TeX root=../main.tex
\begin{table}[!htb]
	\begin{adjustwidth}{-2cm}{0cm}
		\caption[Comparativo]{Comparativo (Paradigma: \ocsex{bol'jiI} `maior')}\label{tab:comparativo}
		\begin{center}
			\begin{tabular}[c]{llllllll}
				\toprule
				                   &                   & \multicolumn{2}{c}{masc.neut. indet.} & \multicolumn{2}{c}{fem. indet.} & \multicolumn{2}{c}{masc.neut. det.}                                                               \\
				\midrule
				\emph{sg.}         & \emph{nom.}       & \cirilex{bol'ijI}                     & \ocsex{bol'ijI}                 & \cirilex{bol'ISi}                   & \ocsex{bol'ija}   & \cirilex{bol'ijI}   & \ocsex{bol'ijI}   \\
				\rowcolor{gray!10} & \emph{acu.}       & \cirilex{bol'ijI}                     & \ocsex{bol'ijI}                 & \cirilex{bol'ISõ}                   & \ocsex{bol'ijõ}   & \cirilex{bol'ijI}   & \ocsex{bol'ijIa}  \\
				                   & \emph{n.-a.neut.} & \cirilex{bol'e}                       & \ocsex{bol'e}                   &                                     &                   & \cirilex{bol'eje}   & \ocsex{bol'eje}   \\
				\rowcolor{gray!10} & \emph{gen.}       & \cirilex{bol'ISa}                     & \ocsex{bol'ISa}                 & \cirilex{bol'ISẽ}                   & \ocsex{bol'ISẽ}   & \cirilex{-Sa(a)go}  & \ocsex{-Sa(a)go}  \\
				                   & \emph{dat.}       & \cirilex{bol'ISu}                     & \ocsex{bol'ISu}                 & \cirilex{bol'ISi}                   & \ocsex{bol'ISi}   & \cirilex{-Su(u)mu}  & \ocsex{-Su(u)mu}  \\
				\rowcolor{gray!10} & \emph{inst.}      & \cirilex{bol'ISemI}                   & \ocsex{bol'ISemI}               & \cirilex{bol'ISejõ}                 & \ocsex{bol'ISejõ} & \cirilex{-Si(ji)mi} & \ocsex{-Si(ji)mi} \\
				                   & \emph{loc.}       & \cirilex{bol'ISi}                     & \ocsex{bol'ISi}                 & \cirilex{bol'ISi}                   & \ocsex{bol'ISi}   & \cirilex{-Si(ji)mi} & \ocsex{-Si(ji)mi} \\
				\midrule
				\emph{du.}         & \emph{nom.acu}    & \cirilex{bol'ISa}                     & \ocsex{bol'ISa}                 & \cirilex{bol'ISi}                   & \ocsex{bol'ISi}   & \cirilex{-Saja}     & \ocsex{-Saja}     \\
				\rowcolor{gray!10} & \emph{n.-a.neut.} & \cirilex{bol'ISi}                     & \ocsex{bol'ISi}                 &                                     &                   & \cirilex{-Siji}     & \ocsex{-Siji}     \\
				                   & \emph{gen.loc.}   & \cirilex{bol'ISu}                     & \ocsex{bol'ISu}                 & \cirilex{bol'ISu}                   & \ocsex{bol'ISu}   & \cirilex{-Suju}     & \ocsex{-Suju}     \\
				\rowcolor{gray!10} & \emph{dat.inst.}  & \cirilex{bol'ISema}                   & \ocsex{bol'ISema}               & \cirilex{bol'ISama}                 & \ocsex{bol'ISama} & \cirilex{-Si(ji)ma} & \ocsex{-Si(ji)ma} \\
				\midrule
				\emph{pl.}         & \emph{nom.}       & \cirilex{bol'ISe}                     & \ocsex{bol'ISe}                 & \cirilex{bol'ISe}                   & \ocsex{bol'ISe}   & \cirilex{-Siji}     & \ocsex{-Siji}     \\
				                   &                   & \cirilex{bol'ISi}                     & \ocsex{bol'ISi}                 &                                     &                   &                     &                   \\
				\rowcolor{gray!10} & \emph{acu.}       & \cirilex{bol'iSẽ}                     & \ocsex{bol'iSẽ}                 & \cirilex{bol'iSẽ}                   & \ocsex{bol'iSẽ}   & \cirilex{-Sẽjẽ}     & \ocsex{-Sẽjẽ}     \\
				                   & \emph{n.-a.neut.} & \cirilex{bol'ISa}                     & \ocsex{bol'ISa}                 &                                     &                   & \cirilex{-Saja}     & \ocsex{-Saja}     \\
				\rowcolor{gray!10} & \emph{gen.}       & \cirilex{bol'ISI}                     & \ocsex{bol'ISI}                 & \cirilex{bol'ISI}                   & \ocsex{bol'ISI}   & \cirilex{-Si(ji)xU} & \ocsex{-Si(ji)xU} \\
				                   & \emph{dat.}       & \cirilex{bol'ISemU}                   & \ocsex{bol'ISemU}               & \cirilex{bol'ISamU}                 & \ocsex{bol'ISamU} & \cirilex{-Si(ji)mU} & \ocsex{-Si(ji)mU} \\
				\rowcolor{gray!10} & \emph{inst.}      & \cirilex{bol'ISi}                     & \ocsex{bol'ISi}                 & \cirilex{bol'ISami}                 & \ocsex{bol'ISami} & \cirilex{-Si(ji)mi} & \ocsex{-Si(ji)mi} \\
				                   & \emph{loc.}       & \cirilex{bol'ISixU}                   & \ocsex{bol'ISixU}               & \cirilex{bol'ISaxU}                 & \ocsex{bol'ISaxU} & \cirilex{-Si(ji)xU} & \ocsex{-Si(ji)xU} \\
				\bottomrule
			\end{tabular}
		\end{center}
	\end{adjustwidth}
\end{table}


Há dois subtipos: o tipo antigo é perpetuado nos comparativos em
-\emph{ijь} (neutro em -\emph{je}), casos oblíquos em -\emph{ьš}-, pouco
numerosos mas que compreendem os adjetivos usuais (\cirilocsex{bol'ijI}
``maior'', \cirilocsex{mIn'ijI} ``menor'', \cirilocsex{luCijI}
``melhor'').
O tipo mais frequente é em -\emph{ějь} (neutro em
-\emph{ěje}), casos oblíquos em -\emph{ějьš-}: \cirilocsex{novEjI},
\gloss{gen.sg.} \cirilocsex{novEjISa}, \gloss{nt.} \cirilocsex{novEje}, \gloss{fem.}
\cirilocsex{novEjISi} ``mais nova''; após consoante palatal, ele apresenta o
alomorfe -\emph{ajь} (neutro \emph{-aje}), casos oblíquos -\emph{ajьš-};
assim, nos temas em velar, a palatalização provocada pelo sufixo de
comparativo faz com que haja uma habitual alternância \emph{k
$\sim$ č}, \emph{g $\sim$ š}, \emph{x
$\sim$ š}, e o \emph{ě} sufixal é despalatalizado em \emph{a}
após a palatal (\emph{cf.}~\ref{sec:estruturasilabica}): \cirilocsex{mUnogU} ``numeroso'',
comparativo \cirilocsex{mUnozajI}, \gloss{gen.} \cirilocsex{mUnozajISa},
\gloss{nt.} \cirilocsex{mUnoZaje}, \gloss{fem.} \cirilocsex{mUnoZajISi}.

O comparativo, por ser um adjetivo, tem uma dupla flexão, determinada e
indeterminada. A flexão indeterminada segue a do tipo em -\emph{jo}-
para os masculinos e neutros (com exceção do nom.acu.sg.), em
-\emph{ja}- para os femininos, a flexão determinada segue a do tipo mole
 \cirilocsex{niStijI}.

Notar-se-á que a oposição entre a forma indeterminada e a forma
determinada é neutralizada nos casos retos do masculino singular
(\cirilocsex{bol'ijI}, \cirilocsex{novEjI} nos dois tipos), no \gloss{loc.pl.}
determinado, a forma contrata \cirilocsex{-ixU} se confunde com o
\gloss{loc.pl.masc.nt.} indeterminado (\cirilocsex{bol'ISixU}), como
sinalizamos acima.

Constata-se no \gloss{nom.pl.masc.} da forma indeterminada uma variação, já que
a forma em \cirilocsex{-ISe}, vestígio do tipo consonantal original,
concorre com uma forma em \cirilocsex{ISi}, analógica do tipo em
-\emph{o}- mole, que segue o alinhamento paradigmático com o tipo
temático. É sobre essa última forma que o \gloss{nom.pl.masc.} determinado em
\cirilocsex{-ISiji} é formado. Encontra-se por vezes no \gloss{acu.sg.masc.} uma
forma \ocsex{bol'ISI}, analógica dos particípios presente e passado
cuja flexão é paralela (\emph{cf.}~\ref{sec:participios}). O \gloss{nom.acu.sg.} neutro determinado,
por sua vez, apresenta duas variantes, a forma etimológica
\cirilocsex{bol'eje} que concorre com uma forma refeita com o sufixo
-\emph{ьš}- do resto da flexão \cirilocsex{bol'ISeje}: é essa
última que é a forma usual em búlgaro antigo e em macedônio antigo, e a
forma em -\emph{eje} é conservada apenas por alguns adjetivos usuais
(\cirilocsex{mIn'eje} ``menor'').

Há, como em várias línguas indo-europeias, comparativos supletivos,
formados a partir de uma raiz diferente do positivo, formando um par
apenas no léxico, não na morfologia: esse é o caso, por exemplo, de
\cirilocsex{luCijI} ``melhor'', que serve de comparativo para \cirilocsex{dobrU}
``bom'', já que seu comparativo etimológico \cirilocsex{dobrEjI} é mais raro;
também é o que ocorre com \cirilocsex{mIn'ijI} ``menor'', que serve de
comparativo a \cirilocsex{malU} ``pequeno''; ou com \cirilocsex{bol'ijI}
``maior'', comparativo de \cirilocsex{velikU} ``grande''.
Os adjetivos em -\emph{kъ}: \cirilocsex{vysokU} ``alto'', \gloss{comparativo.}
\cirilocsex{vySijI}, \cirilocsex{SirokU} ``largo'', \gloss{compar.}
\cirilocsex{Sir'ijI}.

O comparativo de inferioridade é de tipo analítico, formado com um
advérbio \cirilocsex{mIn'e} ``menos'' que precede o adjetivo no positivo
(\cirilocsex{mIn'e dobrU} ``menos bom/pior''). O advérbio \cirilex{mIn'e} é na
verdade o neutro do comparativo \cirilocsex{mIn'ijI} ``menos'', na
forma indeterminada.

O superlativo relativo não tem expressão específica e é a forma
determinada do comparativo que é utilizada nessa função. Às vezes ela é
reforçada por um prefixo \cirilocsex{naji-} (\cirilocsex{najiskorEje}
``o mais rápido'').
Esse procedimento de prefixação é regular para a
expressão do superlativo absoluto, que é formado com o prefixo
\cirilocsex{prE-} (\cirilocsex{prEvelikU} ``muito grande'', forma determinada
\cirilocsex{prEvelikyjI} ``o maior'', \cirilocsex{prEmõdrU} ``muito
sábio'', forma determinada \cirilocsex{prEmõdryjI} ``o mais sábio'');
essa expressão sintética coexiste com uma expressão analítica formada
com o advérbio \cirilocsex{dz'Elo,dzElo} ``muito'' seguido do adjetivo no
positivo (\cirilocsex{dz'Elo velikU,dzElo velikU} ``muito grande'').
\clearpage

\subsection{Particípios}\label{sec:participios}

Os particípios se declinam como adjetivos e têm, como eles, a dupla
flexão determinada e indeterminada. O particípio presente ativo é
formado com o sufixo -\emph{ǫšt-} / \emph{-ęšt-}, que continua o sufixo
indo-europeu *-\emph{nt}- recaracterizado em *-\emph{i̯o}
(*\emph{-ont-jo-}) exceto no Nom.sg.masc.; a alternância \emph{ǫ
$\sim$ ę} não resulta diretamente apenas da harmonia
intrassilábica e sofreu refações. O particípio passado ativo é formado
com o sufixo \emph{-ъs̆-} (< *-\emph{us}-, antigo sufixo do
particípio perfeito indo-europeu), com o alomorfe -\emph{vъs̆-} para os
temas vocálicos. Eles estão, portanto, originalmente, temas
consonantais. O particípio presente passivo é formado com o sufixo
-\emph{mo}-, e o particípio passado passivo com os sufixos -\emph{to}- e
\emph{-no}-; seguem, portanto, a flexão dos temas em -\emph{o}-. O
particípio perfeito é formado com o sufixo -\emph{lo}-, e é, na prática,
defectivo, porque é sempre empregue em uma forma perifrástica com o
auxiliar ``ser'', e tem, portanto, apenas um nominativo (em todos os
gêneros e números), enquanto as outras formas casuais desapareceram já
desde o eslavo comum.

O particípio presente ativo (\autoref{tab:partatvindet}) se divide em três subtipos, um
tipo em -\emph{ǫs̆t-} no nom.sg.masc, em -\emph{y} para os presentes em
-\emph{e}- e os atemáticos, um tipo em -\emph{ęs̆t-} no nom.sg.masc., em
-\emph{ę} para os presentes em -\emph{i}- e um tipo em -\emph{ǫs̆t}- no
nom.sg. em -\emph{ę} para os presentes em -\emph{je}-. A alternância
entre -\emph{y} e -\emph{ę} no nom.sg.masc. é a mesma que a que se
encontra no ac.pl.masc. dos temas em -\emph{o}- e -\emph{jo}-
(\cirilocsex{grady} ``cidades'' < \mbox{*-\emph{nos}}, mas \cirilocsex{kon'ẽ}
``cavalos'' < *-\emph{(j)ens} < *-\emph{(j)nos}, o
primeiro desnasalizado enquanto o segundo permanece nasalizado). Posto à
parte o \gloss{nom.sg.masc.} e neutro, a flexão de três tipos é idêntica.
Notar-se-á uma irregularidade, que vale também para o particípio passado
ativo: a flexão indeterminada tem uma oposição no neutro entre o \gloss{nom.sg.}
e o \gloss{ac.sg.}, ainda que o que defina o gênero neutro seja a ausência de
distinção entre \gloss{nom.} e \gloss{ac.} Isso é consequência de o
\gloss{nom.sg.nt.} e o \gloss{nom.sg.masc.} serem idênticos (neutralização da
oposição de gênero).

% TeX root=../main.tex

\begin{table}[!htb]
	\begin{adjustwidth}{0cm}{0cm}
		\caption[Particípio presente ativo, flexão indeterminada]{Particípio presente ativo, flexão indeterminada}\label{tab:partatvindet}
		\begin{center}
			\begin{tabular}[c]{llllllll}
				\toprule
				                   &                & \multicolumn{2}{c}{masc. indet.}      & \multicolumn{2}{c}{neut. indet.}    & \multicolumn{2}{c}{fem. indet.}                                                         \\
				\midrule
				\emph{sg.}         & \emph{nom.}    & \cirilex{nesy}                        & \ocsex{nesy}                        & \cirilex{nesy}                  & \ocsex{nesy}    & \cirilex{nesõSti} & \ocsex{nesõSti} \\
				\rowcolor{gray!10} & \emph{acu.}    & \cirilex{nesõStI}                     & \ocsex{nesõStI}                     & \cirilex{nesõSte}               & \ocsex{nesõSte} & \cirilex{nesõStõ} & \ocsex{nesõStõ} \\
				                   & \emph{gen.}    & \multicolumn{2}{c}{\cirilex{nesõSta}} & \multicolumn{2}{c}{\ocsex{nesõSta}} & \cirilex{nesõStẽ}               & \ocsex{nesõStẽ}                                       \\
				\midrule
				\emph{du.}         & \emph{nom.acu} & \cirilex{nesõSta}                     & \ocsex{nesõSta}                     & \cirilex{nesõSti}               & \ocsex{nesõSti} & \cirilex{nesõSti} & \ocsex{nesõSti} \\
				\midrule
				\emph{pl.}         & \emph{nom.}    & \cirilex{nesõSte}                     & \ocsex{nesõSte}                     & \cirilex{nesõSta}               & \ocsex{nesõSta} & \cirilex{nesõStẽ} & \ocsex{nesõStẽ} \\
				\rowcolor{gray!10} & \emph{acu.}    & \cirilex{nesõStẽ}                     & \ocsex{nesõStẽ}                     & \cirilex{nesõSta}               & \ocsex{nesõSta} & \cirilex{nesõStẽ} & \ocsex{nesõStẽ} \\
				                   & \emph{gen.}    & \multicolumn{2}{c}{\cirilex{nesõStI}} & \multicolumn{2}{c}{\ocsex{nesõStI}} & \cirilex{nesõStI}               & \ocsex{nesõStI}                                       \\
				\bottomrule
			\end{tabular}
		\end{center}
	\end{adjustwidth}
\end{table}



Seguem o mesmo paradigma os dois outros tipos de particípio presente:
\cirilocsex{molẽ}, \gloss{gen.} \cirilocsex{molẽSta} ``rezante/rezando'' e
\cirilocsex{glagol'ẽ}, \gloss{gen.} \cirilocsex{glagol'õSta} ``falante/falando'',
\cirilocsex{znajẽ}, \gloss{gen.} \cirilocsex{znajõSta} ``sabente/sabendo''.

% TeX root=../main.tex

\begin{table}[!htb]
	\begin{adjustwidth}{0cm}{0cm}
		\caption[Particípio presente ativo, flexão determinada]{Particípio presente ativo, flexão determinada}\label{tab:partatvdet}
		\begin{center}
			\begin{tabular}[c]{llllllll}
				\toprule
				                   &                & \multicolumn{2}{c}{masc. det.}              & \multicolumn{2}{c}{neut. det.}            & \multicolumn{2}{c}{fem. det.}                                                                   \\
				\midrule
				\emph{sg.}         & \emph{nom.}    & \cirilex{nesyjI}                            & \ocsex{nesyjI}                            & \cirilex{nesõSteje}           & \ocsex{nesõSteje}     & \cirilex{nesõStija} & \ocsex{nesõStija} \\
				\rowcolor{gray!10} & \emph{acu.}    & \cirilex{nesõStijI}                         & \ocsex{nesõStijI}                         & \cirilex{nesõSteje}           & \ocsex{nesõSteje}     & \cirilex{nesõStõjõ} & \ocsex{nesõStõjõ} \\
				                   & \emph{gen.}    & \multicolumn{2}{c}{\cirilex{nesõSta(a)go}}  & \multicolumn{2}{c}{\ocsex{nesõSta(a)go}}  & \cirilex{nesõStiji}           & \ocsex{nesõStiji}                                               \\
				\midrule
				\emph{du.}         & \emph{nom.acu} & \cirilex{nesõStaja}                         & \ocsex{nesõStaja}                         & \cirilex{nesõStiji}           & \ocsex{nesõStiji}     & \cirilex{nesõStiji} & \ocsex{nesõStiji} \\
				\midrule
				\emph{pl.}         & \emph{nom.}    & \cirilex{nesõStiji}                         & \ocsex{nesõStiji}                         & \cirilex{nesõSta}             & \ocsex{nesõStaja}     & \cirilex{nesõStẽjẽ} & \ocsex{nesõStẽjẽ} \\
				                   &                & \cirilex{nesõSteji}                         & \ocsex{nesõSteji}                         &                               &                       &                     &                   \\
				\rowcolor{gray!10} & \emph{acu.}    & \cirilex{nesõStẽjẽ}                         & \ocsex{nesõStẽjẽ}                         & \cirilex{nesõStaja}           & \ocsex{nesõStaja}     & \cirilex{nesõStẽjẽ} & \ocsex{nesõStẽjẽ} \\
				                   & \emph{gen.}    & \multicolumn{2}{c}{\cirilex{nesõSti(ji)xU}} & \multicolumn{2}{c}{\ocsex{nesõSti(ji)xU}} & \cirilex{nesõSti(ji)xU}       & \ocsex{nesõSti(ji)xU}                                           \\
				\bottomrule
			\end{tabular}
		\end{center}
	\end{adjustwidth}
\end{table}



A flexão determinada (\autoref{tab:partatvdet}) segue o tipo \cirilocsex{niStijI}.
O \gloss{nom.pl.masc.} em \cirilocsex{-Steji}, com a desinência de tipo consonantal
-\emph{e}, sofre concorrência da forma em -\cirilocsex{Stiji} analógica
dos adjetivos do tipo \cirilocsex{niStiji}.
A irregularidade constituída
pela distinção entre \gloss{nom.sg.} e \gloss{ac.sg.} para o neutro indeterminado não
subsiste na flexão determinada.
A contração de \emph{nesyjь} em
\emph{nesy} neutraliza a oposição entre forma determinada e forma
indeterminada no \gloss{nom.sg.masc.} e transforma várias formas do paradigma em
homófonas (assim o \gloss{nom.pl.masc.} -\emph{s̆tiji > -s̆ti}
homófono do \gloss{nom.sg.fem.} \gloss{indet.} \emph{-s̆ti}).

% TeX root=../main.tex

\begin{table}[!htb]
	\begin{adjustwidth}{0cm}{0cm}
		\caption[Particípio passado passivo, flexão indeterminada]{Particípio passado passivo, flexão indeterminada}\label{tab:partpassindet}
		\begin{center}
			\begin{tabular}[c]{llllllll}
				\toprule
				                   &                & \multicolumn{2}{c}{masc. indet.}     & \multicolumn{2}{c}{neut. indet.}   & \multicolumn{2}{c}{fem. indet.}                                                      \\
				\midrule
				\emph{sg.}         & \emph{nom.}    & \cirilex{nesU}                       & \ocsex{nesU}                       & \cirilex{nesU}                  & \ocsex{nesUSI} & \cirilex{nesISi} & \ocsex{nesISi} \\
				\rowcolor{gray!10} & \emph{acu.}    & \cirilex{nesISI}                     & \ocsex{nesISI}                     & \cirilex{nesISe}                & \ocsex{nesISe} & \cirilex{nesISõ} & \ocsex{nesISõ} \\
				                   & \emph{gen.}    & \multicolumn{2}{c}{\cirilex{nesISa}} & \multicolumn{2}{c}{\ocsex{nesISa}} & \cirilex{nesISẽ}                & \ocsex{nesISẽ}                                     \\
				\midrule
				\emph{du.}         & \emph{nom.acu} & \cirilex{nesISa}                     & \ocsex{nesISa}                     & \cirilex{nesISi}                & \ocsex{nesISi} & \cirilex{nesISi} & \ocsex{nesISi} \\
				\midrule
				\emph{pl.}         & \emph{nom.}    & \cirilex{nesISe}                     & \ocsex{nesISe}                     & \cirilex{nesISa}                & \ocsex{nesISa} & \cirilex{nesISẽ} & \ocsex{nesISẽ} \\
				\rowcolor{gray!10} & \emph{acu.}    & \cirilex{nesISẽ}                     & \ocsex{nesISẽ}                     & \cirilex{nesISa}                & \ocsex{nesISa} & \cirilex{nesISẽ} & \ocsex{nesISẽ} \\
				                   & \emph{gen.}    & \multicolumn{2}{c}{\cirilex{nesISI}} & \multicolumn{2}{c}{\ocsex{nesISI}} & \cirilex{nesISI}                & \ocsex{nesISI}                                     \\
				\bottomrule
			\end{tabular}
		\end{center}
	\end{adjustwidth}
\end{table}


A flexão do particípio passado ativo (\autoref{tab:partpassindet}) é
rigorosamente paralela à do particípio presente.
O mesmo ocorre com os subtipos
\cirilocsex{mol'I}, \gloss{gen.} \cirilocsex{mol'ISa} ``havendo rezado'' (verbos em
-\emph{i}-), e \cirilocsex{byvU}, \gloss{gen.} \cirilocsex{byvUSa} ``tendo sido''
(sobre a formação do tema, \emph{cf.}~\ref{sec:partpass}).

A flexão determinada é também paralela à do particípio presente
(\gloss{nom.sg.masc.} \cirilocsex{nesyjI}, \gloss{acc.} \cirilocsex{nesUSijI},
\gloss{gen.} \cirilocsex{nesUSa(a)go}, etc.; \cirilocsex{mol'ijI}, \gloss{ac.}
\gloss{mol'ISijI}, \gloss{gen.} \cirilocsex{mol'ISa(a)go}, etc.).
O \gloss{nom.acc.sg.nt.} é \cirilocsex{nesUSeje}, \cirilocsex{molISeje}.
A única forma digna de nota é a do \gloss{nom.sg.masc.} \cirilocsex{mol'ijI},
que pode figurar
sob a variante \cirilocsex{mol'ejI} com resultado ocidental do \emph{jer}
``tenso'' (de *\emph{mol'ьjь}, \emph{cf.}~\ref{sec:jers}).
Notar-se-á que para o tipo
\cirilocsex{nesyjI} a forma determinada do particípio passado é idêntica à
do particípio presente no \gloss{nom.sg.masc}.

O particípio presente passivo, que segue o tipo em -\emph{o}- duro, é em
-\emph{mъ}: \cirilocsex{znajemU} ``conhecido'', \cirilocsex{nesomU}
``levado'', \cirilocsex{molimU} ``rezado'', na flexão determinada
\cirilocsex{znajemyjь}, \cirilocsex{nesomyjь} etc.

O particípio passado passivo segue o mesmo tipo em -\emph{o}- duro:
\cirilocsex{nesenU} ``que foi levado'', forma determinada \cirilocsex{nesenyjI},
etc.

O particípio perfeito é defectivo: existe apenas na forma indeterminada
e somente no nominativo nos três números, e não é empregue senão como
núcleo de um tempo perifrástico, com exceção de algumas formas
lexicalizadas que não mais funcionam como particípios. As formas de
nominativo seguem o tipo duro dos temas em -\emph{o}- e -\emph{a}-:
\gloss{sg.masc.} \cirilocsex{-lU}, \gloss{fem.} \cirilocsex{-la},
\gloss{neutro} \cirilocsex{-lo}; \gloss{dual masc.} \cirilocsex{-la},
\gloss{fem.} \cirilocsex{-lE}, neutro \cirilocsex{-lE}; \gloss{pl.masc.}
\cirilocsex{-li}, \gloss{fem.} \cirilocsex{-ly}, \gloss{neutro} \cirilocsex{-la}
(\cirilocsex{znalI jestU} ``ele tem conhecido/conheceu'', \cirilocsex{molila
jestU} ``ela tem rezado/rezou'', \cirilocsex{nesli sõtU} ``eles tem
levado/levaram'').


\section{Pronomes e determinativos}

Há em AEE, como nas outras línguas indo-europeias, dois tipos de
pronomes, os pronomes pessoais e os pronomes não pessoais. As duas
classes têm morfologia diferente que reflete suas diferentes funções.

\subsection{Pronomes pessoais}

As categorias pertinentes são pessoa (P1, P2), número (singular, dual e
plural) e caso. A categoria do gênero não é pertinente, ``eu'' e ``tu''
têm a mesma forma para um referente masculino ou feminino. O número para
os pronomes pessoais não é idêntico ao número para os substantivos,
``nós'' não é o plural de ``eu'', ou seja ``eu'' + ``eu'', mas sim
``eu''+ ``tu''/ ``ele(s)''/ ``ela(s)'' / ``vós''. Há, portanto, três
temas diversos, singular, dual e plural, mas a língua em larga medida os
unificou até fazer com que essencialmente o tema do dual e do plural
coincidissem.

O eslavo herdou do indo-europeu duas grandes oposições na morfologia dos
pronomes pessoais, que são preservadas em AEE. A primeira é uma oposição
entre uma série ``forte'' tônica e uma série ``fraca'' átona, enclítica.
As duas séries se baseiam em temas diversos. Entretanto, como a série
átona é defectiva, a oposição apenas subsiste em AEE no acusativo e
dativo. A série tônica é o termo marcado: emprega-se as formas tônicas
em caso de ênfase (a posição em início de proposição era uma forma de
ênfase) e as formas átonas nos outros casos. Pode-se encontrar outras
formas átonas de acc. Em emprego preposicional (\cirilocsex{na mẽ}
``quanto a mim'', \cirilocsex{na sẽ} ``quanto a si mesmo''), que é
vedado em outras línguas indo-europeias, como o grego, mas jamais se
encontra formas átonas do \gloss{dat.} em emprego preposicional (\emph{*kъ
mi}). A segunda oposição é uma oposição casual, que distingue o
nominativo dos outros casos. O nominativo dos pronomes pessoais não tem
a oposição entre tônico e átono e tem apenas forma tônica, o que é
herdado do indo-europeu: em eslavo, como nas outras línguas
indo-europeias antigas, não há pronome sujeito conjunto; a desinência
verbal é suficiente para exprimir a pessoa, então todo pronome pessoal
sujeito é necessariamente redundante (em relação à desinência verbal), e
ocupa sintaticamente a posição de aposição ao sujeito não explícito,
como o ``moi'' do francês ``\emph{moi, je pense}''. O nominativo é
baseado em um tema diferente do tema tônico do resto da flexão e do tema
átono.

Por fim, existe uma forma específica de pronome pessoal reflexivo, cuja
flexão é paralela à do pronome P2 e conserva características herdadas do
indo-europeu; ele é indiferente ao gênero (como os pronomes pessoais),
ao número (é empregue para se referir a sujeitos singulares, duais e
plurais) e à pessoa (é empregue para se referir a um sujeito que pode
ser tanto P1 como P2 ou um outro sujeito além das pessoas da
interlocução) e é defectivo (não tem forma de nominativo). Ele também
apresenta uma oposição entre formas tônicas e átonas. Para o emprego do
reflexivo, \emph{cf.}~\ref{5.3.2}.

% TeX root=../main.tex
\begin{table}[!h]
	\begin{adjustwidth}{0cm}{0cm}
		\caption[Pronomes pessoais]{Pronomes pessoais}\label{tab:pronomespessoais}
		\begin{center}
			\begin{tabular}[c]{lllllllll}
				\toprule
				                                   & \multicolumn{4}{c}{P1 ``eu''\slash{}``nós'} & \multicolumn{4}{c}{P2 ``tu''\slash{}``vós'}                                                                                                                                    \\
				                                   & \multicolumn{2}{c}{tônico}                  & \multicolumn{2}{c}{átono}                   & \multicolumn{2}{c}{tônico} & \multicolumn{2}{c}{átono}                                                                           \\
				\midrule
				\emph{sg.nom.}                     & \cirilex{azU/jazU}                          & \ocsex{azU/jazU}                            & \cirilex{}                 & \ocsex{}                  & \cirilex{ty}     & \ocsex{ty}         & \cirilex{}      & \ocsex{}      \\
				\rowcolor{gray!10} \emph{acu.}     & \cirilex{mene}                              & \ocsex{mene}                                & \cirilex{mẽ}               & \ocsex{mẽ}                & \cirilex{tebe}   & \ocsex{tebe}       & \cirilex{tẽ}    & \ocsex{tẽ}    \\
				\emph{gen.}                        & \cirilex{mene}                              & \ocsex{mene}                                & \cirilex{}                 & \ocsex{}                  & \cirilex{tebe}   & \ocsex{tebe}       & \cirilex{}      & \ocsex{}      \\
				\rowcolor{gray!10} \emph{dat.}     & \cirilex{mInE}                              & \ocsex{mInE}                                & \cirilex{mi}               & \ocsex{mi}                & \cirilex{tebE}   & \ocsex{tebE}       & \cirilex{ti}    & \ocsex{ti}    \\
				\emph{inst.}                       & \cirilex{mInojõ}                            & \ocsex{mInojõ}                              & \cirilex{}                 & \ocsex{}                  & \cirilex{tobojõ} & \ocsex{tobojõ}     & \cirilex{}      & \ocsex{}      \\
				\rowcolor{gray!10} \emph{loc.}     & \cirilex{mInE}                              & \ocsex{mInE}                                & \cirilex{}                 & \ocsex{}                  & \cirilex{tebE}   & \ocsex{tebE}       & \cirilex{}      & \ocsex{}      \\
				\midrule
				\emph{du.nom.}                     & \cirilex{vE}                                & \ocsex{vE}                                  & \cirilex{}                 & \ocsex{}                  & \cirilex{va}     & \ocsex{va} (bulg.) & \cirilex{}      & \ocsex{}      \\
				                                   &                                             &                                             & \cirilex{}                 & \ocsex{}                  & \cirilex{vy}     & \ocsex{vy} (mac.)  & \cirilex{}      & \ocsex{}      \\
				\rowcolor{gray!10}\emph{acc.}      & \cirilex{na}                                & \ocsex{na} (bulg.)                          & \cirilex{}                 & \ocsex{}                  & \cirilex{va}     & \ocsex{va} (bulg.) & \cirilex{}      & \ocsex{}      \\
				\rowcolor{gray!10}                 & \cirilex{ny}                                & \ocsex{ny} (mac.)                           &                            &                           & \cirilex{vy}     & \ocsex{vy} (mac.)  & \cirilex{}      & \ocsex{}      \\
				\emph{gen.loc}                     & \cirilex{naju}                              & \ocsex{naju}                                & \cirilex{}                 & \ocsex{}                  & \cirilex{vaju}   & \ocsex{vaju}       & \cirilex{}      & \ocsex{}      \\
				\rowcolor{gray!10} \emph{dat.inst} & \cirilex{nama}                              & \ocsex{nama}                                & d.~\cirilex{na}            & d.~\ocsex{na}             & \cirilex{vama}   & \ocsex{vama}       & d.~\cirilex{va} & d.~\ocsex{va} \\
				\midrule
				\emph{pl.nom.}                     & \cirilex{my}                                & \ocsex{my}                                  & \cirilex{}                 & \ocsex{}                  & \cirilex{vy}     & \ocsex{vy}         & \cirilex{}      & \ocsex{}      \\
				\rowcolor{gray!10} \emph{acu.}     & \cirilex{nasU}                              & \ocsex{nasU}                                & \cirilex{ny}               & \ocsex{ny}                & \cirilex{vasU}   & \ocsex{vasU}       & \cirilex{vy}    & \ocsex{vy}    \\
				\emph{gen.}                        & \cirilex{nasU}                              & \ocsex{nasU}                                & \cirilex{}                 & \ocsex{}                  & \cirilex{vasU}   & \ocsex{vasU}       & \cirilex{}      & \ocsex{}      \\
				\rowcolor{gray!10} \emph{dat.}     & \cirilex{namU}                              & \ocsex{namU}                                & \cirilex{nU}               & \ocsex{nU}                & \cirilex{vamU}   & \ocsex{vamU}       & \cirilex{vy}    & \ocsex{vy}    \\
				\emph{inst.}                       & \cirilex{nami}                              & \ocsex{nami}                                & \cirilex{}                 & \ocsex{}                  & \cirilex{vami}   & \ocsex{vami}       & \cirilex{}      & \ocsex{}      \\
				\rowcolor{gray!10} \emph{loc.}     & \cirilex{nasU}                              & \ocsex{nasU}                                & \cirilex{}                 & \ocsex{}                  & \cirilex{vasU}   & \ocsex{vasU}       & \cirilex{}      & \ocsex{}      \\
				\bottomrule
			\end{tabular}
		\end{center}
	\end{adjustwidth}
\end{table}


Paradigma do pronome reflexivo ``se'', que existe apenas no singular:
\gloss{ac.} \cirilocsex{sebe} (átono \cirilocsex{sẽ}), \gloss{gen.}
\cirilocsex{sebe}, \gloss{dat.} \cirilocsex{sebE} (átono \cirilocsex{si}),
\gloss{inst.} \cirilocsex{sobojǫ}, \gloss{loc.} \cirilocsex{sebE}.
É o acusativo átono \cirilocsex{sẽ} que é utilizado para
formar os verbos reflexivos, que são um um dos meios de exprimir o
passivo.

No dual, constata-se uma flutuação dialetal: há apenas uma série e não
duas, com exceção do dativo que opõe uma forma átona \cirilocsex{na},
\cirilocsex{va} à forma tônica de \gloss{dat.inst.} \cirilocsex{nama},
\cirilocsex{vama},
e apenas o búlgaro antigo preserva as formas antigas \cirilocsex{na} e
\cirilocsex{va} no acusativo; em macedônico antigo elas são substituídas
por formas do plural \cirilocsex{ny}, \cirilocsex{vy}, que têm a
particularidade de ser representadas nas duas séries, tônica (\gloss{nom.pl.}
\cirilocsex{vy}) e átona (\gloss{acc.pl.} \cirilocsex{vy}), de onde surge uma
neutralização que permitiu sua extensão ao dual tanto em função do \gloss{nom.}
como do \gloss{ac.} (mas segundo alguns autores \emph{vy} é a forma antiga e
\emph{va} é uma inovação do búlgaro antigo). O tema do dual é o mesmo
que o do plural (exceto no nominativo). Notar-se-á que no plural a série
átona é, na verdade, reduzida a apenas um elemento, e que não há mais
oposição entre \gloss{ac.} e \gloss{dat.} contrariamente ao singular, ao que concorre
para a P2 a homofonia entre a forma átona e a forma do nominativo,
tônica.

Para os pronomes pessoais a forma de \gloss{ac.} é sistematicamente idêntica à
forma de \gloss{gen.}, tanto no singular como no plural (mas não no dual): a
identidade entre \gloss{ac.} e \gloss{gen.} é uma marca de animacidade. Dado que
pronomes pessoais são por definição animados e determinados, e são mais
elevados que os substantivos em uma escala de animacidade em todas as
línguas que apresentam esse critério, também não seria surpreendente que
o \gloss{gen.-acc.} animado deles seja menos limitado e tenha uma extensão maior
(não se limita ao singular).

Não há pronome de terceira pessoa, é o anafórico, que não é um pronome
pessoal, que é empregue nessa função.

\subsection{Adjetivos possessivos}

Aos pronomes pessoais correspondem os adjetivos possessivos (\autoref{tab:adjposs}),
determinantes que funcionam também como pronomes. Eles seguem a flexão
do tipo pronominal mole (representado pelo demonstrativo \cirilocsex{sь}
``este'', \emph{cf.}~\autoref{tab:prondem}). Indicam o gênero do possuído, mas não o do
possessor, em razão da indistinção de gênero dos pronomes pessoais. Há
também um adjetivo possessivo reflexivo, que segue o mesmo tipo
flexional e que é, assim como o pronome reflexivo, indiferente ao
número; mas não há possessivo para a terceira pessoa não reflexivo:
emprega-se nessa função o genitivo do anafórico (\cirilocsex{jego}
``dele'', \cirilocsex{jejẽ} ``dela'', \cirilocsex{ixъ} ``deles/delas'', com
marcação do gênero do possessor no singular mas não no dual ou plural, e
sem marcação de gênero do possuído para os três números). Essa
distribuição é herdada do indo-europeu. Para o emprego do possessivo
reflexivo, \emph{cf.}~\ref{5.3.2}.

\clearpage%

% TeX root=../main.tex
\begin{table}[!h]
	\begin{adjustwidth}{0cm}{0cm}
		\caption[Adjetivos possessivos]{Adjetivos possessivos}\label{tab:adjposs}
		\begin{center}
			\begin{tabular}[c]{lllllllll}
				\toprule
				                                     & \multicolumn{2}{c}{``meu''}          & \multicolumn{2}{c}{``minha''}      & \multicolumn{2}{c}{``nosso''}        & \multicolumn{2}{c}{``nossa''}                                                                              \\
				                                     & \multicolumn{2}{c}{masc.nt.}         & \multicolumn{2}{c}{fem.}           & \multicolumn{2}{c}{masc.nt.}         & \multicolumn{2}{c}{fem.}                                                                                   \\
				\midrule
				\emph{sg.nom.}                       & \cirilex{mojI}                       & \ocsex{mojI}                       & \cirilex{moja}                       & \ocsex{moja}                       & \cirilex{naSI}   & \ocsex{naSI}   & \cirilex{naSa}   & \ocsex{naSa}   \\
				\rowcolor{gray!10} \emph{acu.}       & \multicolumn{2}{c}{= nom. ou acu.}   & \cirilex{mojõ}                     & \ocsex{mojõ}                         & \multicolumn{2}{c}{= nom. ou acu.} & \cirilex{naSõ}   & \ocsex{naSõ}                                       \\
				\emph{nom.acu.nt.}                   & \cirilex{moje}                       & \ocsex{moje}                       & \cirilex{}                           & \ocsex{}                           & \cirilex{naSe}   & \ocsex{naSe}   & \cirilex{}       & \ocsex{}       \\
				\rowcolor{gray!10} \emph{gen.}       & \cirilex{mojego}                     & \ocsex{mojego}                     & \cirilex{mojejẽ}                     & \ocsex{mojejẽ}                     & \cirilex{naSego} & \ocsex{naSego} & \cirilex{naSejẽ} & \ocsex{naSejẽ} \\
				\emph{dat.}                          & \cirilex{mojemu}                     & \ocsex{mojemu}                     & \cirilex{mojeji}                     & \ocsex{mojeji}                     & \cirilex{naSemu} & \ocsex{naSemu} & \cirilex{naSeji} & \ocsex{naSeji} \\
				\rowcolor{gray!10} \emph{inst.}      & \cirilex{mojimI}                     & \ocsex{mojimI}                     & \cirilex{mojejõ}                     & \ocsex{mojejõ}                     & \cirilex{naSimI} & \ocsex{naSimI} & \cirilex{naSejõ} & \ocsex{naSejõ} \\
				\emph{loc.}                          & \cirilex{mojemI}                     & \ocsex{mojemI}                     & \cirilex{mojeji}                     & \ocsex{mojeji}                     & \cirilex{naSemI} & \ocsex{naSemI} & \cirilex{naSeji} & \ocsex{naSeji} \\
				\midrule
				\emph{du.nom.acu.}                   & \cirilex{moja}                       & \ocsex{moja}                       & \cirilex{moji}                       & \ocsex{moji}                       & \cirilex{naSa}   & \ocsex{naSa}   & \cirilex{naSi}   & \ocsex{naSi}   \\
				\rowcolor{gray!10}\emph{nom.acu.nt.} & \cirilex{moji}                       & \ocsex{moji}                       & \cirilex{}                           & \ocsex{}                           & \cirilex{naSi}   & \ocsex{naSi}   & \cirilex{}       & \ocsex{}       \\
				\emph{gen.loc}                       & \multicolumn{2}{c}{\cirilex{mojeju}} & \multicolumn{2}{c}{\ocsex{mojeju}} & \multicolumn{2}{c}{\cirilex{naSeju}} & \multicolumn{2}{c}{\ocsex{naSeju}}                                                                         \\
				\rowcolor{gray!10}\emph{dat.inst.}   & \multicolumn{2}{c}{\cirilex{mojima}} & \multicolumn{2}{c}{\ocsex{mojima}} & \multicolumn{2}{c}{\cirilex{naSima}} & \multicolumn{2}{c}{\ocsex{naSima}}                                                                         \\
				\midrule
				\emph{pl.nom.}                       & \cirilex{moji}                       & \ocsex{moji}                       & \cirilex{mojẽ}                       & \ocsex{mojẽ}                       & \cirilex{naSi}   & \ocsex{naSi}   & \cirilex{naSẽ}   & \ocsex{naSẽ}   \\
				\rowcolor{gray!10} \emph{acu.}       & \cirilex{mojẽ}                       & \ocsex{mojẽ}                       & \cirilex{mojẽ}                       & \ocsex{mojẽ}                       & \cirilex{naSẽ}   & \ocsex{naSẽ}   & \cirilex{naSẽ}   & \ocsex{naSẽ}   \\
				\emph{nom.acu.nt.}                   & \cirilex{moja}                       & \ocsex{moja}                       & \cirilex{}                           & \ocsex{}                           & \cirilex{naSa}   & \ocsex{naSa}   & \cirilex{}       & \ocsex{}       \\
				\rowcolor{gray!10}\emph{gen.}        & \multicolumn{2}{c}{\cirilex{mojixU}} & \multicolumn{2}{c}{\ocsex{mojixU}} & \multicolumn{2}{c}{\cirilex{naSixU}} & \multicolumn{2}{c}{\ocsex{naSixU}}                                                                         \\
				\emph{dat.}                          & \multicolumn{2}{c}{\cirilex{mojimU}} & \multicolumn{2}{c}{\ocsex{mojimU}} & \multicolumn{2}{c}{\cirilex{naSimU}} & \multicolumn{2}{c}{\ocsex{naSimU}}                                                                         \\
				\rowcolor{gray!10}\emph{inst.}       & \multicolumn{2}{c}{\cirilex{mojimi}} & \multicolumn{2}{c}{\ocsex{mojimi}} & \multicolumn{2}{c}{\cirilex{naSimi}} & \multicolumn{2}{c}{\ocsex{naSimi}}                                                                         \\
				\emph{loc.}                          & \multicolumn{2}{c}{\cirilex{mojixU}} & \multicolumn{2}{c}{\ocsex{mojixU}} & \multicolumn{2}{c}{\cirilex{naSixU}} & \multicolumn{2}{c}{\ocsex{naSixU}}                                                                         \\
				\bottomrule
			\end{tabular}
		\end{center}
	\end{adjustwidth}
\end{table}


Os possessivos \cirilocsex{tvojь} ``teu'' e \cirilocsex{svojь} ``seu
próprio'' (possessivo reflexivo) se declinam como \cirilocsex{mojь}
``meu''; \cirilocsex{vaSь} ``vosso'' se declina como \cirilocsex{naSь}
``nosso''.

\subsection{Pronomes não-pessoais}

Os pronomes não pessoais são os demonstrativos,
interrogativos-indefinidos, relativos e seus derivados. Eles distinguem
as categorias de gênero, número e caso, como as formas nominais. Como as
formas nominais, eles têm o subgênero animado (genitivo-acusativo)
apenas no singular. Como os paradigmas nominais, distingue-se o tipo
duro e o mole. Características da flexão pronominal são a neutralização
da oposição de gênero nos casos oblíquos do plural e a homofonia do
gen.pl. e do loc. Pl. o que se encontra na flexão determinada do
adjetivo.

\clearpage%
\subsubsection{Demonstrativos}

% TeX root=../main.tex
\begin{table}[!h]
	\begin{adjustwidth}{0cm}{0cm}
		\caption[Pronomes demonstrativos]{Pronomes demonstrativos (Paradigmas: tipo duro \emph{tъ} ``esse'', tipo mole \emph{sь} ``este'')}\label{tab:prondem}
		\begin{center}
			\begin{tabular}[c]{lllllllll}
				\toprule
				                                     & \multicolumn{2}{c}{masc.nt. duro}         & \multicolumn{2}{c}{fem. duro}    & \multicolumn{2}{c}{masc.nt. mole}  & \multicolumn{2}{c}{fem. mole}                                                                                         \\
				\midrule
				\emph{sg.nom.}                       & \cirilex{tU}                              & \ocsex{tU}                       & \cirilex{ta}                       & \ocsex{ta}                                & \cirilex{sI}         & \ocsex{sI}         & \cirilex{si}   & \ocsex{si}   \\
				                                     &                                           &                                  &                                    &                                           & \cirilex{sijI, sejI} & \ocsex{sijI, sejI} &                &              \\
				\rowcolor{gray!10} \emph{acu.}       & \multicolumn{2}{c}{\emph{= nom. ou acu.}} & \cirilex{tõ}                     & \ocsex{tõ}                         & \multicolumn{2}{c}{\emph{= nom. ou acu.}} & \cirilex{sijõ}       & \ocsex{sijõ}                                       \\
				\emph{nom.acu.nt.}                   & \cirilex{to}                              & \ocsex{to}                       & \cirilex{}                         & \ocsex{}                                  & \cirilex{se}         & \ocsex{se}         & \cirilex{}     & \ocsex{}     \\
				\rowcolor{gray!10} \emph{gen.}       & \cirilex{togo}                            & \ocsex{togo}                     & \cirilex{tojẽ}                     & \ocsex{tojẽ}                              & \cirilex{sego}       & \ocsex{sego}       & \cirilex{sejẽ} & \ocsex{sejẽ} \\
				\emph{dat.}                          & \cirilex{tomu}                            & \ocsex{tomu}                     & \cirilex{toji}                     & \ocsex{toji}                              & \cirilex{semu}       & \ocsex{semu}       & \cirilex{seji} & \ocsex{seji} \\
				\rowcolor{gray!10} \emph{inst.}      & \cirilex{tEmI}                            & \ocsex{tEmI}                     & \cirilex{tojõ}                     & \ocsex{tojõ}                              & \cirilex{simI}       & \ocsex{simI}       & \cirilex{sejõ} & \ocsex{sejõ} \\
				\emph{loc.}                          & \cirilex{tomI}                            & \ocsex{tomI}                     & \cirilex{toji}                     & \ocsex{toji}                              & \cirilex{semI}       & \ocsex{semI}       & \cirilex{seji} & \ocsex{seji} \\
				\midrule
				\emph{du.nom.acu.}                   & \cirilex{ta}                              & \ocsex{ta}                       & \cirilex{tE}                       & \ocsex{tE}                                & \cirilex{sija}       & \ocsex{sija}       & \cirilex{si}   & \ocsex{si}   \\
				\rowcolor{gray!10}\emph{nom.acu.nt.} & \cirilex{tE}                              & \ocsex{tE}                       & \cirilex{}                         & \ocsex{}                                  & \cirilex{si}         & \ocsex{si}         & \cirilex{}     & \ocsex{}     \\
				\emph{gen.loc}                       & \multicolumn{2}{c}{\cirilex{toju}}        & \multicolumn{2}{c}{\ocsex{toju}} & \multicolumn{2}{c}{\cirilex{seju}} & \multicolumn{2}{c}{\ocsex{seju}}                                                                                      \\
				\rowcolor{gray!10}\emph{dat.inst.}   & \multicolumn{2}{c}{\cirilex{tEma}}        & \multicolumn{2}{c}{\ocsex{tEma}} & \multicolumn{2}{c}{\cirilex{sima}} & \multicolumn{2}{c}{\ocsex{sima}}                                                                                      \\
				\midrule
				\emph{pl.nom.}                       & \cirilex{ti}                              & \ocsex{ti}                       & \cirilex{ty}                       & \ocsex{ty}                                & \cirilex{siji, si}   & \ocsex{siji, si}   & \cirilex{sijẽ} & \ocsex{sijẽ} \\
				\rowcolor{gray!10} \emph{acu.}       & \cirilex{ty}                              & \ocsex{ty}                       & \cirilex{ty}                       & \ocsex{ty}                                & \cirilex{sijẽ}       & \ocsex{sijẽ}       & \cirilex{sijẽ} & \ocsex{sijẽ} \\
				\emph{nom.acu.nt.}                   & \cirilex{ta}                              & \ocsex{ta}                       & \cirilex{}                         & \ocsex{}                                  & \cirilex{si}         & \ocsex{si}         & \cirilex{}     & \ocsex{}     \\
				\rowcolor{gray!10}\emph{gen.}        & \multicolumn{2}{c}{\cirilex{tExU}}        & \multicolumn{2}{c}{\ocsex{tExU}} & \multicolumn{2}{c}{\cirilex{sixU}} & \multicolumn{2}{c}{\ocsex{sixU}}                                                                                      \\
				\emph{dat.}                          & \multicolumn{2}{c}{\cirilex{tEmU}}        & \multicolumn{2}{c}{\ocsex{tEmU}} & \multicolumn{2}{c}{\cirilex{simU}} & \multicolumn{2}{c}{\ocsex{simU}}                                                                                      \\
				\rowcolor{gray!10}\emph{inst.}       & \multicolumn{2}{c}{\cirilex{tEmi}}        & \multicolumn{2}{c}{\ocsex{tEmi}} & \multicolumn{2}{c}{\cirilex{simi}} & \multicolumn{2}{c}{\ocsex{simi}}                                                                                      \\
				\emph{loc.}                          & \multicolumn{2}{c}{\cirilex{tExU}}        & \multicolumn{2}{c}{\ocsex{tExU}} & \multicolumn{2}{c}{\cirilex{sixU}} & \multicolumn{2}{c}{\ocsex{sixU}}                                                                                      \\
				\bottomrule
			\end{tabular}
		\end{center}
	\end{adjustwidth}
\end{table}


O AEE tem três demonstrativos (\autoref{tab:prondem}), \cirilocsex{sь} ``este'' é o
dêitico proximal (<\emph{*k̑i-}), \cirilocsex{onъ} ``esse'' (que
segue a flexão de \emph{tъ}) o dêitico distal
(*\emph{h\textsubscript{1}e/ono-}), \cirilocsex{tъ} o dêitico ``central''
não marcado (< *\emph{to-}). O valor pessoal dos dêiticos associa
\emph{sь} à primeira pessoa, \emph{tъ} à segunda e \emph{onъ} à
terceira, mas isso é raramente perceptível em AEE exceto quando o
contexto opõe os temas pronominais um ao outro. Há um quarto termo,
\cirilocsex{ovU} (que segue a flexão de \emph{tъ}), que é empregue quase
exclusivamente nas oposições binárias \cirilocsex{ovU... ovU}
``esse\ldots{} aquele''.

Entre o tipo duro e o tipo brando há as mesmas alomorfias que nos temas
nominais, \emph{ě $\sim$ i} (\gloss{Inst.sg.}, \gloss{Gen.dat.inst.loc.pl.},
\gloss{n.acc.du.fem.} e \gloss{nt.}), \emph{o $\sim$ e}
(\gloss{gen.dat.loc.sg.masc.}, \gloss{n.acc.sg.nt.}), \emph{ъ $\sim$ ь}
(\gloss{nom.sg.masc.}).
A flexão de \cirilocsex{sь} não segue exatamente a flexão
do tipo pronominal brando (representado pelo possessivo \cirilocsex{naSь}
``nosso'') mas comporta um certo número de formas específicas, devidas
ao fato que o tema alterna entre \emph{se-} e \emph{si(j)-} (de fato
\emph{sьj-} com \emph{jer} tenso, \emph{cf.}~\ref{sec:jers}).
No \gloss{nom.sg.masc.}, as formas
\cirilocsex{sijь} e \cirilocsex{sejь} são recentes, elas figuram sobretudo no
\glsxtrshort{Supr}.

Seguem a flexão de \cirilocsex{tъ} os pronomes e determinantes
\cirilocsex{onъ} ``aquele'', \cirilocsex{ovъ} ``o outro (de dois)'',
\cirilocsex{samъ} ``(ele) mesmo'', \cirilocsex{inъ} ``outro'', a série
\cirilocsex{takъ} ``tal'', \cirilocsex{kakъ} ``(de) qual (tipo)'',
\cirilocsex{nEkakъ} ``(de) algum (tipo)'', \cirilocsex{nikakъ} ``(de) nenhum
(tipo)'', \cirilocsex{jakъ(Ze)} ``tal que'', \cirilocsex{inakъ} ``(de)
outro (tipo)'', \cirilocsex{vъsEkъ} ``(de) toda (sorte)'' (com a aplicação
geral da alternância dos temas em velar; \gloss{loc.sg.masc.} \cirilocsex{tacE},
\gloss{nom.pl.masc.} \cirilocsex{taci}, etc.), assim como os cardinais
\cirilocsex{jedinъ} ``um'' e \cirilocsex{dъva} ``dois'', este último não tem
formas de dual.

% TeX root=../main.tex
\begin{table}[!h]
	\begin{adjustwidth}{0cm}{0cm}
		\caption[Paradigma misto de \emph{vьsь} ``todo'', \emph{sicь} ``(de) certo (tipo)'']{Paradigma misto de \emph{vьsь} ``todo'', \emph{sicь} ``(de) certo (tipo)''}\label{tab:pronmist}
		\begin{center}
			\begin{tabular}[c]{lllllll}
				\toprule
				                                & \multicolumn{2}{c}{masc.nt.}              & \multicolumn{2}{c}{fem.}           & \multicolumn{2}{c}{masc.nt.}                                                                                 \\
				\midrule
				\emph{sg.nom.}                  & \cirilex{vIsI}                            & \ocsex{vIsI}                       & \cirilex{vIsE}               & \ocsex{vIsE}                              & \cirilex{sicI}   & \ocsex{sicI}   \\
				                                &                                           &                                    & \cirilex{vIsa}               & \ocsex{vIsa}                              &                  &                \\
				\rowcolor{gray!10} \emph{acu.}  & \multicolumn{2}{c}{\emph{= nom. ou acu.}} & \cirilex{vIsjõ}                    & \ocsex{vIs'õ}                & \multicolumn{2}{c}{\emph{= nom. ou acu.}}                                     \\
				\rowcolor{gray!10}              & \multicolumn{2}{c}{                     } & \cirilex{vIsõ}                     & \ocsex{vIsõ}                 & \multicolumn{2}{c}{                     }                                     \\
				\emph{nom.acu.nt.}              & \cirilex{vIse}                            & \ocsex{vIse}                       & \cirilex{}                   & \ocsex{}                                  & \cirilex{sice}   & \ocsex{sice}   \\
				\rowcolor{gray!10} \emph{gen.}  & \cirilex{vIsego}                          & \ocsex{vIsego}                     & \cirilex{vIsejẽ}             & \ocsex{vIsejẽ}                            & \cirilex{sicego} & \ocsex{sicego} \\
				\emph{dat.}                     & \cirilex{vIsemu}                          & \ocsex{vIsemu}                     & \cirilex{vIseji}             & \ocsex{vIseji}                            & \cirilex{sicemu} & \ocsex{sicemu} \\
				\rowcolor{gray!10} \emph{inst.} & \cirilex{vIsEmI}                          & \ocsex{vIsEmI}                     & \cirilex{vIsejõ}             & \ocsex{vIsejõ}                            & \cirilex{sicEmI} & \ocsex{sicEmI} \\
				\emph{loc.}                     & \cirilex{vIsemI}                          & \ocsex{vIsemI}                     & \cirilex{vIseji}             & \ocsex{vIseji}                            & \cirilex{sicemI} & \ocsex{sicemI} \\
				\midrule
				\emph{pl.nom.}                  & \cirilex{vIsi}                            & \ocsex{vIsi}                       & \cirilex{vIsẽ}               & \ocsex{vIsẽ}                              & \cirilex{sici}   & \ocsex{sici}   \\
				\rowcolor{gray!10} \emph{acu.}  & \cirilex{vIsẽ}                            & \ocsex{vIsẽ}                       & \cirilex{vIsẽ}               & \ocsex{vIsẽ}                              & \cirilex{sicẽ}   & \ocsex{sicẽ}   \\
				\emph{nom.acu.nt.}              & \cirilex{vIsE}                            & \ocsex{vIsE}                       & \cirilex{}                   & \ocsex{}                                  & \cirilex{sica}   & \ocsex{sica}   \\
				                                & \cirilex{vIsa}                            & \ocsex{vIsa}                       & \cirilex{}                   & \ocsex{}                                  & \cirilex{}       & \ocsex{}       \\
				\rowcolor{gray!10}\emph{gen.}   & \multicolumn{2}{c}{\cirilex{vIsExU}}      & \multicolumn{2}{c}{\ocsex{vIsExU}} & \cirilex{sicExU}             & \ocsex{sicExU}                                                                \\
				\emph{dat.}                     & \multicolumn{2}{c}{\cirilex{vIsEmU}}      & \multicolumn{2}{c}{\ocsex{vIsEmU}} & \cirilex{sicEmU}             & \ocsex{sicEmU}                                                                \\
				\rowcolor{gray!10}\emph{inst.}  & \multicolumn{2}{c}{\cirilex{vIsEmi}}      & \multicolumn{2}{c}{\ocsex{vIsEmi}} & \cirilex{sicEmi}             & \ocsex{sicEmi}                                                                \\
				\emph{loc.}                     & \multicolumn{2}{c}{\cirilex{vIsExU}}      & \multicolumn{2}{c}{\ocsex{vIsExU}} & \cirilex{sicExU}             & \ocsex{sicExU}                                                                \\
				\bottomrule
			\end{tabular}
		\end{center}
	\end{adjustwidth}
\end{table}


Dois pronomes, o indefinido \cirilocsex{vьsь} ``todo'' e o demonstrativo
\cirilocsex{sicь} ``d(este) tipo'', têm um paradigma irregular misturando
o tipo brando e o tipo duro (tabela 24). O paradigma original é de tipo
duro (conservado no \gloss{inst.sg.masc.nt.} e nos casos oblíquos do plural),
mas por causa da terceira palatalização a maior parte das formas casuais
foram passados ao tipo brando. O dual não é atestado para
\cirilocsex{vьsь} e é suplementado pela forma \cirilocsex{oba} ``todos os
dois'' (\gloss{fem.} e \gloss{nt.} \cirilocsex{obE}, que segue a flexão de dual de
\cirilocsex{dъva} ``dois''). No nom.sg.fem. e no nom.acc.pl.nt regularmente
ocorre \cirilocsex{sica}, mas para \emph{vьsь} coexistem duas grafias,
\cirilocsex{vьsě} e \cirilocsex{vsja,vs'a} ou \cirilocsex{vьsa}, que são resultado
da dificuldade de grafar o /s'/ palatalizado (cf.3.1.1.). A grafia com
-\emph{a} é regular nos manuscritos búlgaros antigos, esporádica nos
manuscritos macedônicos antigos que conservam o uso do glagolítico para
notar \emph{ě}. Paralelamente, a grafia \cirilex{vIsõ} em vez de
\cirilocsex{vIsjõ,vьs'ǫ}
para o \gloss{acc.sg.fem.} é regular nos manuscritos búlgaros antigos, mas também
em boa parte dos manuscritos em glagolítico, enquanto a grafia \cirilex{vIsjõ} não é
encontrada senão nos manuscritos mais recentes (\emph{Ps., Euch.}). No
dialeto morávio, a forma é \emph{vьšь} (o tratamento *\emph{x}
> \emph{š} para a segunda e a terceira palatalizações é
característico do eslavo ocidental).


\subsubsection{Interrogativo--indefinido}

O interrogativo-indefinido (tabela 25) continua as duas funções herdadas
do indo-europeu (tema *\emph{kwi-}, *\emph{kwe/o}-). O pronome
interrogativo-indefinido só existe no singular. O tema é \emph{ko-} para
os animados (sem distinção de gênero entre masculino e feminino, e com o
\gloss{acc}. Idêntico ao gen. característico dos animados),
\emph{čь-}/\emph{če-} para o inanimado, com a forma de \gloss{nom.sg.}
\cirilocsex{kъto}, \cirilocsex{Cьto} recaracterizada.

Tabela 25: Pronome interrogativo-indefinido

% TODO: TABELA

A flexão do inanimado é dupla, com uma forma ``longa'' na qual a forma
do genitivo \cirilocsex{Ceso} (< *\emph{kwso}, gr. τέο/τοῦ) cumpre
a função de tema ao qual são acrescidas as desinências do animado, e uma
forma ``curta'', paralela à do animado, tomando como tema \cirilocsex{Ce-}
(variante \cirilocsex{Cь-} por analogia do nom.acc.); no dativo as formas
\cirilocsex{Cesomu}, \cirilocsex{Cьsomu} são regulares em búlgaro antigo
como em macedônio antigo (o dativo não é atestado no \emph{Missal de
Kiev}), e a forma curta \cirilocsex{Cmeu} é excepcional. A forma negativa
é \cirilocsex{nikъto(Ze)} ``ninguém'', \cirilocsex{niCьtoZe} ``nada'',
reforçada pela partícula \cirilocsex{že}, e com as mesmas variantes
(\emph{ničьso} no genitivo no \emph{Missal de Kiev}): em emprego
preposicional a preposição se insere imediatamente diante do indefinido,
após a negação \cirilocsex{ni-} (\cirilocsex{ni o komь Ze} ``sobre
ninguém''). Encontra-se excepcionalmente uma forma de inanimado
\cirilocsex{niCьZe} ``nada'' (\emph{niCьZe} \emph{Ps.} 38,6). A partir de
\cirilex{kU-to} é formado o indefinido \cirilocsex{kъZьdo} ``cada um'', que conserva a
forma antiga do nom.sg. \cirilocsex{kъ} (*\emph{kwos}) não recaracterizada
por \cirilocsex{-to}.

As formas simples \cirilocsex{kъto}, \cirilocsex{čьto} são preservadas em
emprego indefinido apenas em proposições negativas (\emph{ni sъmě kъto
vъprositi jego} ``e ninguém ousou lhe perguntar'' Mt. 22, 46),
interrogativas (\emph{imate li čьto}? ``tendes algo?'' L 24,41) ou
hipotéticas (\emph{ašte kъto}... ``se alguém...''), assim como em
algumas proposições subordinadas subjuntivas; em outros contextos o
indefinido foi recaracterizado e substituído pela forma composta
\cirilocsex{nEkъto} ``alguém'', \cirilocsex{nECto} ``algo'', o que lembra a
distribuição latina entre \emph{quis} e \emph{aliquis}.

A partir do interrogativo é formado o adijetivo \cirilocsex{Cijь} ``de
quem?'', que segue a flexão de \cirilocsex{mojь} (nom.sg.fem.
\cirilocsex{Cija} ``de quem?'', com possuído feminino, nom.sg.nt.
\cirilocsex{Cije}).

O AEE distingue o pronome interrogativo-indefinido \cirilocsex{kъto},
\cirilocsex{Cьto} do determinante interrogativo-indefinido \cirilocsex{kyjь}
``qual'' (*\emph{kъ-jь}, com dois elementos flexionados originalmente, o
que ainda é visível nas formas de casos retos no singular e plural (daí
a palatalização no nom.pl.masc. \cirilocsex{ci-ji}), mas não nos casos
oblíquos (tabela 26). Como ocorre com o pronome, o AEE reserva
\cirilocsex{kyjь} para o emprego interrogativo e utiliza, para o emprego
indefinido, a forma recaracterizada \cirilocsex{nEkyjь} ``algum''. A forma
negativa é \cirilocsex{nikyjьZe} ``nenhum''. Como ocorre com o pronome, a
preposição se intercala entre o negativo e o indefinido (\emph{ni otъ
kojejęže strasti} ``de mal algum'' \glsxtrshort{Supr}). A partir de
\cirilocsex{kyjь} é formado o adjetivo indefinido \cirilocsex{kyjьZьdo}
``cada'', que corresponde ao pronome \cirilocsex{kъZьdo}. Um outro
interrogativo-indefinido é \cirilocsex{kotoryjь} ``qual (entre muitos)?'',
que segue a flexão regular de \cirilocsex{novyjь}; \cirilocsex{kotoryjь} é
às vezes pronome outras determinante, daí seus compostos
\cirilocsex{nEkotoryjь} ``algum'', \cirilocsex{nikotoryjьZe}
``nenhum''.

Tabela 26: Determinante interrogativo-indefinido \emph{kyjь}

% TODO: TABELA

Uma única forma de dual é atestada, o \gloss{dat.inst.} \cirilocsex{kyjima}.

\subsubsection{Anafórico e relativo}\label{sec:anafóricoerelativo}

O pronome anafórico (Tabela 27) é baseado em um tema alternante
\emph{i-} (<\emph{\#ji-}) $\sim$ \emph{je-}; nessa
tabela, conserva-se a notação \emph{\#jь} e não \emph{\#i}, que é a
transliteração usual, porque os dois são fonologicamente distintos (cf.
3.2.3.), em contrapartida, nota-se \emph{\#i-} em lugar de *\emph{ji-},
porque os dois se confundem: todas as formas em \emph{i-} resultam de um
\emph{*ji}- mais antigo, que se encontra em posição intervocálica na
flexão determinada do adjetivo (\emph{cf.}~\ref{sec:adjparadigma}). 
O anafórico é
defectivo e não há forma de nominativo, é o nominativo do demonstrativo
\cirilocsex{tъ} que é utilizada em função de anafórico por supletismo (entre
parênteses na tabela 27).
O acc.sg.masc. não tem forma de animado de data antiga:
nos \emph{Evangelhos}, o acc. É regularmente \cirilocsex{jь},
(transliterado \emph{i}) tanto para animados como para inanimados
(\emph{iscěli i} ``ele o cura'' L 22, 51). O alinhamento com os outros
pronomes, contudo, suscitou a substituição de /\cirilocsex{jь} por
\cirilocsex{jego}, forma do genitivo, já nos \emph{Evangelhos}
(\emph{prědajęi ego} ``aquele que o trai'' M 14,55), portanto, com um
genitivo-acusativo. A forma \emph{jego} tende, de fato, a se estender
igualmente para os inanimados, e o mesmo vale para o relativo
\cirilocsex{iZe} (\emph{kamenь jegože běšę privęzali} ``a pedra que eles
firmaram'' \emph{Supr.} 154, 17). No plural, o acc. \cirilocsex{jẽ} começa
a sofrer concorrência do gen.-acc. \cirilocsex{ixъ}, seguindo o modelo do
singular (\emph{vъprosi ję} ``ele os interrogou'' Mt. 22, 41
\emph{Mar.}, mas \emph{Sav.} tem \emph{ixъ}).

O pronome relativo, por sua vez, é derivado do anafórico por acréscimo
da partícula \cirilocsex{Ze} e segue, portanto, a mesma flexão:
regularmente tem forma de nominativo. Historicamente, os dois pronomes,
relativo e anafórico, são derivados do mesmo tema pronominal *\emph{jo-}
que se produziu o relativo grego e indo-iraniano: sob a nova forma
\emph{iže}, ele conservou sua função relativa, enquanto a forma antiga
se tornava anafórica. Esse é o pronome que, cliticizado, fornece a
flexão determinada de adjetivo em eslavo e báltico.

Tabela 27: Anafórico-relativo

% TODO: TABELA

O anafórico e o relativo têm a particularidade de apresentar, após
preposição, um alomorfe com nasal inicial: \cirilocsex{otъ n'ego}
``dele'', \cirilocsex{prědъ n'imi} ``diante deles'', \cirilocsex{o
n'eji} ``a seu sujeito'' (fem.), relativo \cirilocsex{na n'emьZe}
``sobre o qual'', \cirilocsex{na n'ejiZe} ``sobre a qual'',
\cirilocsex{iz n'egoZe} ``do qual''; para o \gloss{acc.sg.masc.} a forma
nasalizada é
\cirilocsex{n'ь} (em comparação com \cirilocsex{jь} para a forma regular).
Essa nasal é, originalmente, a consoante final de algumas preposições
(\cirilocsex{vъ(n)} ``em'', \cirilocsex{sъ(n)} ``com''), emudecida em razão do
emudecimento das consoantes finais, mas conservada diante de anafórico,
em razão do caráter proclítico da preposição, e reanalisado como uma
marca de anafórico ({[}vъn'emь{]} ``nele'', morfologicamente \emph{vъn
jemь}, silabificado {[}vъ\textbar n'e\textbar mь{]} e, por conseguinte,
reanalisado como \emph{vъ n'emь}. O alomorfe com nasal inicial foi, em
seguida, generalizado após todas as preposições, inclusive aquelas que
não tinham nasal final originalmente. O anafórico e o relativo,
portanto, apresentam uma distribuição complementar entre um alomorfe com
nasal prefixada, em emprego preposicional, e um alomorfe sem nasal, em
todos os outros contextos.


\section{Numerais}

Os quatro primeiros cardinais são determinantes que concordam em gênero
e caso com o substantivo. O cardinal \cirilocsex{jedinъ} ``um'' se
declina como o demonstrativo \cirilocsex{tъ}. O cardinal \cirilocsex{dъva} tem
uma flexão de dual regular (\gloss{nom.acc.masc.} \cirilocsex{dъva}, \gloss{nt.}
\cirilocsex{dъvE}, \gloss{fem.} \cirilocsex{dъvE}, \gloss{gen.loc.}
\cirilocsex{dъvoju} para os três
gêneros, \gloss{dat.inst.} \cirilocsex{dъvEma} para os três gêneros); 
a flexão é a mesma para \cirilocsex{oba} ``ambos os dois''.
Os cardinais \cirilocsex{trije} ``três'' e \cirilocsex{Cetyre} ``quatro'' são
\emph{pluralia tantum} que seguem a flexão dos temas em \emph{-i-} para
``três''
(\gloss{nom.pl.masc.} \cirilocsex{trije}, \gloss{nom.pl.fem.} e \gloss{nt.}
\cirilocsex{tri}, \gloss{gen.pl.} \cirilocsex{trijь} para os três gêneros) e dos
temas consonantais para ``quatro'' (\gloss{nom.pl.masc.} \cirilocsex{Cetyre,}
\gloss{nom.pl.fem.} e \gloss{nt.} \cirilocsex{Cetyri}, \gloss{gen.pl.}
\cirilocsex{Cetyrъ} para os três gêneros, com uma variante
\cirilocsex{Cetyrь}).

Os cardinais de ``cinco'' a ``nove'' (cf. Anexo 1) são substantivos, que
são coletivos femininos em \emph{-i-} e seguem, portanto, a flexão de
\cirilocsex{kostь}: embora originalmente coletivos, não têm plural nem
dual. Como são substantivos, eles são seguidos de complemento nominal de
genitivo plural (\cirilocsex{devẽtь mõZь} ``nove homens'', \emph{lit.}
``uma novena de homens'').

O cardinal \cirilocsex{desẽtь} ``dez'' segue a flexão do tipo consonantal
e se declina como \cirilocsex{dьnь}, com exceção do \gloss{inst.sg.} que é
\cirilocsex{desẽtijõ} (ou \cirilex{desẽtIjõ}), como o feminino \cirilocsex{kostijõ}. De
fato, a analogia dos cardinais precedentes fez com que houvesse a
tendência à passagem ao tipo em -\emph{i}-. Ele também é um substantivo
com complemento no genitivo plural. O mesmo vale para \cirilocsex{sъto}
``cem'', neutro em -\emph{o}-, \cirilocsex{tysõSti} ``mil'', feminino em
-\emph{i}- e \cirilocsex{tьma} ``dez mil'', feminino em \emph{-a-}.

Os cardinais de ``onze'' a ``dezenove'' são sintagmas compostos da
unidade (no caso desejado) seguido de \cirilocsex{na desẽte} ``sobre
dez'' (\cirilocsex{dъva na desẽte} ``doze'', \emph{lit.} ``dois
sobre dez'', \gloss{loc.} \cirilocsex{na dъvoju na desẽte} ``sobre doze'').

As dezenas são formadas com o cardinal seguido do numeral ``dez'' no
caso desejado (\cirilocsex{dъva desẽti} ``vinte'', \emph{lit.} ``duas
dezenas'', no dual \cirilocsex{trije desẽte} ``trinta'', no
plural, \cirilocsex{pẽtь desẽtъ} ``cinquenta'', no genitivo plural). Em
``vinte'', ``trinta'' e ``quarenta'', os dois elementos são declinados,
a partir de ``cinquenta'', apenas o multiplicador é declinado, enquanto
o nome da dezena permanece no genitivo plural. Como ``dez'', esses
numerais se constroem com um complemento no genitivo plural (cf.
5.3.5.).

As dezenas e as unidades são coordenadas por \emph{i} ``e'' (por
\emph{ti} no \glsxtrshort{Supr}): \cirilocsex{Cetyre desẽte i dъva}
``quarenta e dois''.

Há algumas formas de coletivo: \cirilocsex{dъvoji} ``grupo de dois'', com
a forma preposicional \cirilocsex{na dъvoje} ``em dois'',
\cirilocsex{troji} ``grupo de três'', \cirilocsex{Cetvorъ} ``grupo de
quatro'', \cirilocsex{desẽtoro} ``grupo de dez''.

As frações são expressas por meio de sintagmas com \cirilocsex{polъ}
``metade'', \cirilocsex{Cetvrьtь} ``um quarto'', \cirilocsex{desẽtina} ``um
décimo'' seguidos de complemento nominal no genitivo.

Os ordinais seguem a flexão regular dos adjetivos, e têm as duas formas,
determinado e indeterminado. São baseados sobretudo sobre o tema do
cardinal com o sufixo *-\emph{to}-, com exceção dos dois primeiros que
são supletivos (cf. Anexo 1). Os ordinais de ``onze'' a ``dezenove'' são
formados com o ordinal da unidade e o cardinal da dezena, que permanece
invariável (\emph{vъtoroe na desęte} ``décimo segundo'' no neutro, lit.
``segundo sobre dez''). Os ordinais das dezenas são formados pelo
acréscimo do sufixo ao cardinal composto (\cirilocsex{dъvadesẽtьnъ}
``vigésimo''), formação bastante recente. Os ordinais compostos são
formados com um ordinal da unidade e um ordinal da dezena, coordenados
por \emph{i} ``e'' ou sem coordenação.

\section{Verbos}
\lipsum[20]

\subsection{Categorias}
\lipsum[19]

\subsubsection{Voz}
\lipsum[18]

\subsubsection{Modo}
\lipsum[17]

\subsubsection{Tempo}
\lipsum[16]

\subsubsection{Aspecto}
\lipsum[15]

\subsubsection{Formas nominais}
\lipsum[14]

\subsection{Tipos flexionais}\label{sec:tiposflexionais}
\lipsum[13]

\subsection{Tema de presente}
\lipsum[12]

\subsubsection{Presente}
\lipsum[11]

\subsubsection{Particípio presente}
\lipsum[10]

\subsubsection{Imperativo}
\lipsum[9]

\subsection{Tema de infinitivo}
\lipsum[8]

\subsubsection{Aoristo}
\lipsum[7]

\subsubsection{Imperfeito}
\lipsum[6]

\subsubsection{Infinitivo e supino}
\lipsum[5]

\subsubsection{Particípios passados}\label{sec:partpass}
\lipsum[4]

\subsection{Formas perifrásticas}
\lipsum[3]

\subsection{Emprego dos temas aspectuais}
\lipsum[2]

\subsection{Emprego dos modos}
\lipsum[1]
